%-----------------------------------------------------------------------------
% Rozdział IV
%-----------------------------------------------------------------------------

\cleardoublepage
\chapter{Gomulscy z~Anieliny, Wólki Mińskiej i~Karoliny}

% Przednia okładka podrozdziału
\includepdf{Anielina_mapa_fin.png}

\phantomsection
\addcontentsline{toc}{section}{Anielina: 1853~r. -- \biezacyrok}

\lettrine[lines=5, lhang=0.1]{A}{nielina} została założona około 1820~roku
w~odległości dwóch kilometrów na~południe od Mińska Mazowieckiego, na terenie
ówczesnego Królestwa Polskiego\footnote{Od 1~stycznia 1986~roku Anielina stała
się częścią Mińska Mazowieckiego i~nie funkcjonuje już jako samodzielna
miejscowość.}. Wieś ta należała na początku XIX~wieku do parafii pod
wezwaniem Narodzenia Najświętszej Maryi Panny w~Mińsku Mazowieckim.
Na~początku swojego istnienia wieś Anielina nosiła następujące nazwy w~księgach
metrykalnych parafii w~Mińsku Mazowieckim: \enquote{Kolonia Anielin},
\enquote{Anielin} oraz \enquote{Angelin}, dopiero z~czasem utrwaliła się
obecnie stosowana nazwa tej miejscowości. Najstarsza znana wzmianka
o~Anielinie pochodzi z~1822~roku -- jest to wpis w~aktach stanu cywilnego gminy
Mińsk Mazowiecki, z~pierwszego lipca ww. roku, dotyczący zgonu Marianny
Zawadzkiej córki Jana Zawadzkiego i~Apolonii z~domu Kossakowskiej (patrz: ryc.
\ref{fig:anielina_1822}). Zmarła wówczas Marianna Zawadzka urodziła się
niespełna klika miesięcy wcześniej w~sąsiedniej wsi~Zakole. Poza wspomnianą
rodziną Zawadzkich\footnote{Co ciekawe Zawadzcy mieszkający we~wsi Anielina
w~latach dwudziestych XIX~wieku to przodkowie Marcina i~Marianny Zawadzkich,
którzy przeprowadzili się na przełomie XIX~i~XX~wieku do nowo powstałej
wówczas wsi Desna (obecnie Desno) -- byli jedną z~rodzinzałożycielskich tej
miejscowości (udziałowcami spółki Desna) -- losy rodziny Zawadzkich spod
Mińska Mazowieckiego zostaną szerzej omowione
w~\hyperref[sec:zawadzcy]{załączniku numer IV}~do niniejszej książki.}, wieś
Anielina na początku XIX~wieku zamieszkiwały również między innymi takie
rodziny pochodzenia chłopskiego jak: Sikorscy, Ostrowscy, Baranowie, Michalscy,
Sadowscy czy Nalazkowie.

\begin{figure}[!ht]
    \vspace*{0.5cm}
    \centering \includegraphics[width=0.8\linewidth]{
        Anielina_1822.jpg}
    \captionsetup{format=hang}
    \caption{Pierwszy znany wpis dotyczący wsi Anielina pochodzący 
    z~1822~roku z~akt stanu cywilnego gminy Mińsk Mazowiecki 
    \cite{par_minsk1}.}
    \label{fig:anielina_1822}
\end{figure}

Pierwszym dzieckiem Piotra\indeksosoby{Piotr Franciszek}{Gumulski}{1815}{Rudzienko}
i~Marianny\indeksosoby[Chłopik]{Marianna}{Gumulska}{1819}{Zamienie} Gumulskich,
które urodziło się w~Anielinie, a~zarazem ich czwartą córką, była
Rozalia\indeksosoby[Gomulska]{Rozalia}{Szostak}{1853}{Anielina} Gomulska.
Urodziła się ona 18~listopada 1853~roku, a~w~akcie jej chrztu jej ojciec
został określony jako gospodarz mający 38~lat zamieszkały we~wsi Angelina.
Rodzicami chrzestnymi Rozalii Gomulskiej zostali Władysław Michalski oraz
Józefa Sadowska (patrz: ryc. \ref{fig:rgomulska_1853}).

\begin{figure}[!ht]
    \vspace*{0.5cm}
    \centering \includegraphics[width=1.0\linewidth]{
        1853_Rozalia_Gomulska_Szostak_akt_chrztu_parafia_Mińsk_Mazowiecki_wpis_260.jpg}
    \captionsetup{format=hang}
    \caption{Akt chrztu Rozalii Gomulskiej -- par. Mińsk Mazowiecki 1851~rok 
    (260/1853) \cite{par_minsk2}.}
    \label{fig:rgomulska_1853}
\end{figure}

Dnia 7~października 1857~roku urodziła się najmłodsza córka Piotra i~Marianny
Gumulskich -- Józefa
Gomulska\indeksosoby{Józefa}{Gomulska}{1857}{Anielina}\footnote{Drzewo
genealogiczne najbliższej rodziny Piotra i~Marianny Gumulskich zostało
zaprezentowane w~\hyperref[sec:drzewo_pfgumulski]{załączniku numer I}~do
niniejszej książki.}. W~akcie jej chrztu, jej ojciec został określony jako
\enquote{... kowal, lat czterdzieści pięć liczący, w~Angelinie zamieszkały...}.
Rodzicami chrzestnymi Józefy Gomulskiej zostali Wojciech Chłopik oraz Marianna
Zawadzka (patrz: ryc. \ref{fig:jgomulska_1857}). We~wrześniu 1859~roku we~wsi
Anielina na~świat przyszedł Michał
Gomulski\indeksosoby{Michał}{Gomulski}{1859}{Anielina}, przedostatni syn
małżeństwa Gumulskich, a~zarazem jeden z~trzech członków rodziny Gomulskich,
którzy byli założycielami wsi Desna (a~zarazem udziałowcami spółki
\enquote{Desna}) -- bezpośredni potomkowie Michała Gomulskiego mieszkają
we~wsi Desno do~dnia dzisiejszego, nosząc nazwisko po~swoim przodku. W~akcie
chrztu Michała jego ojciec ponownie został określony jako kowal z~Angeliny.
Rodzicami chrzestnymi Michała Gomulskiego zostali Jan Ostrowski oraz Anna
Wójcik (patrz: ryc. \ref{fig:mgomulski_1859}). 

\newpage
\begin{figure}[!ht]
    \vspace*{0.5cm}
    \centering \includegraphics[width=1.0\linewidth]{
        1857_Józefa_Gomulska_akt_chrztu_parafia_Mińsk_Mazowiecki_wpis_191.jpg}
    \captionsetup{format=hang}
    \caption{Akt chrztu Józefy Gomulskiej -- par. Mińsk Mazowiecki 1857~rok 
    (191/1857) \cite{par_minsk2}.}
    \label{fig:jgomulska_1857}
\end{figure}

Ostatnim dzieckiem Piotra i~Marianny Gumulskich był
Stanisław\indeksosoby{Stanisław}{Gomulski}{1863}{Anielina} urodzony
w~Anielinie 2~listopada 1863~roku. W~momencie jego narodzin jego rodzice mieli
odpowiednio 48~lat (ojciec) oraz 44~lata (matka). Ojciec Stanisława w~akcie
jego chrztu został określony jako kolonista zamieszkały 
w~Angelinie\footnote{W tamtym czasie kolonistami nazywani byli nowo przybyli 
do danej miejscowości osadnicy.}. Rodzicami chrzestnymi Stanisława Gomulskiego 
zostali Władysław Michalski oraz Katarzyna Bodzionek (patrz: ryc. 
\ref{fig:sgomulski_1863}). Stanisław, podobnie jak jego cztery lata starszy 
brat Michał oraz dwa lata starszy bratanek był jedynym z~założycieli wsi
Desna -- bezpośredni potomkowie Stanisława Gomulskiego mieszkają we~wsi Desno
do~dnia dzisiejszego, ale nie noszą oni już nazwiska swojego przodka. Piotr
i~Marianna Gumulscy mieli łącznie jedenaścioro dzieci, z~czego tylko sześcioro
dożyło dorosłości, byli nimi:
Jan\indeksosoby{Jan}{Gomulski}{1836}{Zamienie} (ur. 1836~r. -- zm. 1901~r.),
Grzegorz\indeksosoby{Grzegorz}{Gomulski}{1840}{Grzebowilk} (ur. 1840~r. --
Anna\indeksosoby[Gomulska]{Anna}{Michalska}{1843}{Grzebowilk} (ur. 1843~r. --
zm. 1905~r.),
Rozalia\indeksosoby[Gomulska]{Rozalia}{Szostak}{1853}{Anielina} (ur. 1853~r.
-- zm. 1929~r.), Michał\indeksosoby{Michał}{Gomulski}{1859}{Anielina} 
(ur. 1859~r. -- zm. 1918~r.)
i~Stanisław\indeksosoby{Stanisław}{Gomulski}{1863}{Anielina} (ur. 1863~r.~--
zm. 1929~r.).

\begin{figure}[!ht]
    \vspace*{0.5cm}
    \centering \includegraphics[width=0.82\linewidth]{
        1859_Michał_Gomulski_akt_chrztu_parafia_Mińsk_Mazowiecki_wpis_241.jpg}
    \captionsetup{format=hang}
    \caption{Akt chrztu Michała Gomulskiego -- par. Mińsk Mazowiecki 1859~rok 
    (241/1859) \cite{par_minsk2}.}
    \label{fig:mgomulski_1859}
\end{figure}

\begin{figure}[!ht]
    \vspace*{0.5cm}
    \centering \includegraphics[width=0.82\linewidth]{
        1863_Stanisław_Gomulski_akt_chrztu_parafia_Mińsk_Mazowiecki_wpis_277.jpg}
    \captionsetup{format=hang}
    \caption{Akt chrztu Stanisława Gomulskiego -- par. Mińsk Mazowiecki 
    1863~rok (277/1863) \cite{par_minsk2}.}
    \label{fig:sgomulski_1863}
\end{figure}

Piotr Franciszek Gumulski\indeksosoby{Piotr
Franciszek}{Gumulski}{1815}{Rudzienko} ostanie 17~lat swojego życia przeżył
we~wsi Anielina -- zmarł 8~lutego 1870~roku w~wieku 55~lat (patrz: ryc.
\ref{fig:pfgomulski_1870}\footnote{Akt zgonu Piotra Franciszka Gumulskiego
został sporządzony w~języku rosyjskim, gdyż począwszy od~1868~roku w~wyniku 
represji carskich po powstaniu styczniowym, wszystkie kościelne akta 
metrykalne w~Królestwie Polskim musiały być prowadzone w~języku rosyjskim -- 
represje te~były utrzymywane aż~do~1915~roku.}). Marianna
Gumulska\indeksosoby[Chłopik]{Marianna}{Gumulska}{1819}{Zamienie}
(z~domu Chłopik) zmarła 10~lat po swoim mężu -- 12~kwietnia 1880~roku,
dożywając wieku 61~lat. W~akcie jej zgonu wskazano, iż~zmarła ona we~wsi
Jarząbek, czyli pobliskiej w~stosunku do Anieliny wsi, do~której Marianna
prawdopodobnie przeprowadziła się ze~swoim szwagrem oraz jego rodziną (patrz:
ryc. \ref{fig:mgomulska_1880}).

\begin{figure}[!ht]
    \vspace*{0.5cm}
    \centering \includegraphics[width=0.94\linewidth]{
        1870_Piotr_Franciszek_Gumulski_akt_zgonu_parafia_Mińsk_Mazowiecki_wpis_51.jpg}
    \captionsetup{format=hang}
    \caption{Akt zgonu Piotra Franciszka Gumulskiego -- par. Mińsk Mazowiecki 
    1870~rok (51/1870) \cite{par_minsk2}.}
    \label{fig:pfgomulski_1870}
\end{figure}

\begin{figure}[!ht]
    \vspace*{0.5cm}
    \centering \includegraphics[width=0.94\linewidth]{
        1880_Marianna_Gumulska_Chłopik_akt_zgonu_parafia_Mińsk_Mazowiecki_wpis_71.jpg}
    \captionsetup{format=hang}
    \caption{Akt zgonu Marianny Gumulskiej -- par. Mińsk Mazowiecki 
    1880~rok (71/1880) \cite{par_minsk2}.}
    \label{fig:mgomulska_1880}
\end{figure}

Około 1872~roku, dwa lata po śmierci Piotra Gumulskiego, do Anieliny 
przeprowadził się jego młodszy brat
Andrzej\indeksosoby{Andrzej}{Gumulski}{1829}{Zamienie}. Powodem jego
przeprowadzki do~Anielny był prawdopodobnie ślub, który wziął z~lokalną
dziewczyną, Magdaleną
Michalską\indeksosoby[Michalska]{Magdalena}{Gumulska}{1848}{Anielina},
22~stycznia 1872~roku w~Mińsku Mazowieckim. Andrzej Gumulski miał wówczas
blisko 42~lata, a~jego przyszła żona ledwo ukończyła 23~lata (patrz: ryc.
\ref{fig:agumulski_1872}). 

\begin{figure}[!ht]
    \vspace*{0.5cm}
    \centering \includegraphics[width=0.72\linewidth]{
        1872_Andrzej_Gomulski_Magdalena_Michalska_akt_ślubu_parafia_Mińsk_Mazowiecki_wpis_3.jpg}
    \captionsetup{format=hang}
    \caption{Akt ślubu Andrzeja Gumulskiego i~Magdaleny Michalskiej -- par. 
    Mińsk Mazowiecki 1872~rok (3/1872) \cite{par_minsk2}.}
    \label{fig:agumulski_1872}
\end{figure}

Andrzej Gumulski wraz ze~swoją żoną Magdaleną dochowali się trójki dzieci, 
z~czego dwójka najstarszych urodziła się w~Anielinie, a~trzecie z~nich 
urodziło się we~wsi Jarząbek\footnote{Drzewo genealogiczne najbliższej rodziny
Andrzeja i~Magdaleny Gumulskich zostało zaprezentowane
w~\hyperref[sec:drzewo_agumulski]{załączniku numer I}~do niniejszej książki.}: 

\begin{itemize}
    \item Feliks Gomulski\indeksosoby{Feliks}{Gomulski}{1875}{Anielina}
    (ur. 1875~r. -- zm. 1941~r.) -- patrz: ryc. \ref{fig:fgomulski_1875},
    \item Antonina
    Gomulska\indeksosoby[Gomulska]{Antonina}{Araźna}{1878}{Anielina}
    (ur. 1878~r. -- zm. 1955~r.) -- patrz: ryc. \ref{fig:agomulska_1878},
    \item Ludwika\indeksosoby{Ludwika}{Gomulska}{1882}{Jarząbek} (ur. 1882~r.
    -- zm. 1884~r.)  -- patrz: ryc. \ref{fig:lgomulska_1882} oraz ryc.
    \ref{fig:lgomulska_1884}.
  \end{itemize}

Wygląda więc na~to, że~Andrzej Gumulski wraz z~żoną, dziećmi oraz 
bratową Marianną przeprowadzili się do~Jarząbka około 1880~roku, jednak
po~1884~roku wrócili do~Anieliny, bo~tam zmarł zarówno
Andrzej\indeksosoby{Andrzej}{Gumulski}{1829}{Zamienie} -- w~1887~roku
(patrz: ryc. \ref{fig:agomulski_1887}), jak i~jego żona
Magdalena\indeksosoby[Michalska]{Magdalena}{Gumulska}{1848}{Anielina} --
w~1904~roku.

 \begin{figure}[!ht]
    \vspace*{0.5cm}
    \centering \includegraphics[width=0.88\linewidth]{
        1875_Feliks_Gomulski_akt_chrztu_parafia_Mińsk_Mazowiecki_wpis_126.jpg}
    \captionsetup{format=hang}
    \caption{Akt chrztu Feliksa Gomulskiego -- par. Mińsk Mazowiecki 
    1875~rok (126/1875) \cite{par_minsk2}.}
    \label{fig:fgomulski_1875}
\end{figure}

\begin{figure}[!ht]
    \vspace*{0.5cm}
    \centering \includegraphics[width=0.88\linewidth]{
        1878_Antonina_Gomulska_Araźna_akt_chrztu_parafia_Mińsk_Mazowiecki_wpis_47.jpg}
    \captionsetup{format=hang}
    \caption{Akt chrztu Antoniny Gomulskiej -- par. Mińsk Mazowiecki 
    1878~rok (47/1878) \cite{par_minsk2}.}
    \label{fig:agomulska_1878}
\end{figure}

\begin{figure}[!ht]
    \vspace*{0.5cm}
    \centering \includegraphics[width=1.0\linewidth]{
        1882_Ludwika_Gomulska_akt_chrztu_parafia_Mińsk_Mazowiecki_wpis_167.jpg}
    \captionsetup{format=hang}
    \caption{Akt chrztu Ludwiki Gomulskiej -- par. Mińsk Mazowiecki 
    1882~rok (167/1882) \cite{par_minsk2}.}
    \label{fig:lgomulska_1882}
\end{figure}

\begin{figure}[!ht]
    \vspace*{1cm}
    \centering \includegraphics[width=1.0\linewidth]{
        1884_Ludwika_Gomulska_akt_zgonu_parafia_Mińsk_Mazowiecki_wpis_70.jpg}
    \captionsetup{format=hang}
    \caption{Akt zgonu Ludwiki Gomulskiej -- par. Mińsk Mazowiecki 
    1884~rok (70/1884) \cite{par_minsk2}.}
    \label{fig:lgomulska_1884}
\end{figure}

\begin{figure}[!ht]
    \vspace*{0.5cm}
    \centering \includegraphics[width=0.9\linewidth]{
        1887_Andrzej_Gomulski_akt_zgonu_parafia_Mińsk_Mazowiecki_wpis_13.jpg}
    \captionsetup{format=hang}
    \caption{Akt zgonu Andrzeja Gumulskiego -- par. Mińsk Mazowiecki 
    1887~rok (13/1887) \cite{par_minsk2}.}
    \label{fig:agomulski_1887}
\end{figure}

Z trójki dzieci Andrzeja i~Magdaleny Gumulskich tylko dwójka dożyła wieku 
dorosłego -- Feliks Gomulski\indeksosoby{Feliks}{Gomulski}{1875}{Anielina}
ożenił się w~1909~roku w~Pęcicach pod Pruszkowem z~Małgorzatą
Makowską\indeksosoby[Makowska]{Małgorzata}{Gomulska}{1888}{Słomin}
(ur.~1888~r.~-- zm.~1951~r.) pochodzącą ze~Słomina (patrz: ryc.
\ref{fig:fgomulski_1909}) a~Antonina
Gomulska\indeksosoby[Gomulska]{Antonina}{Araźna}{1878}{Anielina} w~1896~roku
wyszła za~mąż w~Mińsku Mazowieckim za~Michała
Araźnego\indeksosoby{Michał}{Araźny}{1869}{Olesin} (ur.~1869~r.~--
zm.~1953~r.) pochodzącego ze~wsi Olesin, znajdującej się koło Dębego Wielkiego.
Antonina i~Michał Araźni po~ślubie zamieszkali w~Anielinie, gdzie doczekali
się co~najmniej szóstki dzieci:
Stanisława\indeksosoby{Stanisław}{Araźny}{1897}{Anielina} (ur.~1897~r.~--
zm.~1989~r.),
Władysławy\indeksosoby[Araźna]{Władysława}{Klimkowska}{1899}{Anielina}
(ur.~1899~r.~-- zm.~1984~r.),
Zofii\indeksosoby{Zofia}{Araźna}{1901}{Anielina} (ur.~1901~r.~--
zm.~1983~r.), Józefa\indeksosoby{Józef}{Araźny}{1903}{Anielina} (ur.~1903~r.~--
zm.~?), Walerii\indeksosoby{Waleria}{Araźna}{1906}{Anielina} (ur.~1906~r.~--
zm.~?), Aleksandra\indeksosoby{Aleksander}{Araźny}{1909}{Anielina}
(ur.~1909~r.~-- zm.~?) oraz
Anieli\indeksosoby{Aniela}{Araźna}{1911}{Anielina} (ur.~1911~r.~-- zm.~?). 
Feliks Gomulski\indeksosoby{Feliks}{Gomulski}{1875}{Anielina} razem ze~swoją
małżonką przez pierwsze kilka lat po~ślubie mieszkali w~rodzinnej wsi
Małgorzaty\indeksosoby[Makowska]{Małgorzata}{Gomulska}{1888}{Słomin} --
w~Słominie, o~czym świadczy fakt, iż~ich dwóch pierwszych synów
Henryk\indeksosoby{Henryk}{Gomulski}{1910}{Słomin} (ur.~1910~r.~-- zm.~?) oraz
Czesław\indeksosoby{Czesław}{Gomulski}{1912}{Słomin} (ur.~1912~r.~--
zm.~1946~r.) urodziło się (odpowiednio w~1910 oraz 1912~roku) właśnie w~tej
miejscowości (patrz: ryc. \ref{fig:hgomulski_1910} oraz
\ref{fig:czgomulski_1912}). Trzeci syn małżeństwa Gomulskich,
Walerian\indeksosoby{Walerian \flagaPL}{Gomulski}{1914}{Mińsk Mazowiecki}
(ur.~1914~r.~-- zm.~1995~r.), urodził się w~1914~roku w~Mińsku Mazowieckim
(patrz: ryc. \ref{fig:wgomulski_1914}), a~ich pierwsza córka,
Helena\indeksosoby{Helena}{Gomulska}{1917}{Anielina} (ur.~1917~r.~--
zm.~1917~r.), urodziła się (patrz: ryc. \ref{fig:hgomulska_1917}) oraz zmarła
(patrz: ryc. \ref{fig:hgomulska_1917_2}) w~1917~roku w~Anielinie. Oznacza~to,
że~po~około ośmiu latach od~ślubu z~Małgorzatą, Feliks Gomulskich powrócił
do~swojej rodzinnej miejscowości lub w~jej okolice, bo~jego kolejne dzieci
rodziły się zarówno w~Anielinie, jak~i~w~Mińsku Mazowieckim.

\begin{figure}[!ht]
    \vspace*{0.5cm}
    \centering \includegraphics[width=1.0\linewidth]{
        1909_Feliks_Gomulski_Małgorzata_Makowska_akt_ślubu_parafia_Pęcice_wpis_16.jpg}
    \captionsetup{format=hang}
    \caption{Akt ślubu Feliksa Gomulskiego oraz Małgorzaty Makowskiej -- par. 
    Pęcice 1909~rok (16/1909) \cite{par_pecice}.}
    \label{fig:fgomulski_1909}
\end{figure}

Czwarty syn małżeństwa Gomulskich,
Zygmunt\indeksosoby{Zygmunt}{Gomulski}{1918}{Anielina} (ur.~1918~r.~--
zm.~1943~r.), urodził się w~Anielinie, w~1918~roku (patrz: ryc.
\ref{fig:zgomulski_1918}), natomiast ich druga córka,
Janina\indeksosoby[Kubiciel]{Janina}{Gomulska}{1920}{Mińsk Mazowiecki},
(ur.~1920~r.~-- zm.~2008~r.) podobnie jak jej starszy brat
Walerian\indeksosoby{Walerian \flagaPL}{Gomulski}{1914}{Mińsk Mazowiecki},
urodziła się w~Mińsku Mazowieckim, w~1920~roku (patrz: ryc.
\ref{fig:jgomulska_1920}). Co~ciekawe, w~aktach chrztu Heleny, Zygmunta oraz
Janiny, ich ojciec, Feliks Gomulski określony został jako kowal, w~dwóch
pierwszych wskazano, iż~zamieszkuje on~w~Anielinie natomiast w~trzecim, jako
jego miejscowość zamieszkania wskazano Grzebowilk. Oznacza~to, iż~Feliks
Gomulski, kontynuował rodzinną tradycję i~podobnie jak jego ojciec i~dziadek
trudnił się pracą kowala w~miejscowościach, w~których zamieszkiwał.

\begin{figure}[!ht]
    \vspace*{0.5cm}
    \centering \includegraphics[width=0.73\linewidth]{
        1910_Henryk_Gomulski_akt_chrztu_parafia_Pęcice_wpis_77.jpg}
    \captionsetup{format=hang}
    \caption{Akt chrztu Henryka Gomulskiego -- par. Pęcice 1910~rok (77/1910) 
    \cite{par_pecice}.}
    \label{fig:hgomulski_1910}
\end{figure}

\begin{figure}[!ht]
    \vspace*{0.5cm}
    \centering \includegraphics[width=0.73\linewidth]{
        1912_Czesław_Gomulski_akt_chrztu_parafia_Pęcice_wpis_16.jpg}
    \captionsetup{format=hang}
    \caption{Akt chrztu Czesława Gomulskiego -- par. Pęcice 1912~rok (16/1912) 
    \cite{par_pecice}.}
    \label{fig:czgomulski_1912}
\end{figure}

\begin{figure}[!ht]
    \vspace*{0.5cm}
    \centering \includegraphics[width=0.78\linewidth]{
        1914_Walerian_Gomulski_akt_chrztu_parafia_Mińsk_Mazowiecki_wpis_674.jpg}
    \captionsetup{format=hang}
    \caption{Akt chrztu Waleriana Gomulskiego -- par. Mińsk Mazowiecki 
    1914~rok (674/1914) 
    \cite{par_minsk2}.}
    \label{fig:wgomulski_1914}
\end{figure}

\begin{figure}[!ht]
    \vspace*{1cm}
    \centering \includegraphics[width=0.78\linewidth]{
        1917_Helena_Gomulska_akt_chrztu_parafia_Mińsk_Mazowiecki_wpis_50.jpg}
    \captionsetup{format=hang}
    \caption{Akt chrztu Heleny Gomulskiej -- par. Mińsk Mazowiecki 
    1917~rok (50/1917) 
    \cite{par_minsk2}.}
    \label{fig:hgomulska_1917}
\end{figure}

Autor niniejszej książki nie dysponuje informacjami na temat miejsca narodzin, 
dwóch najmłodszych dzieli Feliksa i~Małgorzaty Gomulskich --
Jadwigi\indeksosoby{Jadwiga}{Gomulska}{1923}{Anielina} (ur.~1923~r.~--
zm.~1994~r.) i~Mieczysława\indeksosoby{Mieczysław}{Gomulski}{1929}{Anielina}
(ur.~1929~r.~-- zm.~1968~r.), gdyż w~momencie pisania niniejszej książki,
kościelne akta chrztów dostępne do publicznego wglądu, kończyły się,
w~zależności od~parafii, na~około 1920~roku\footnote{Większość archidiecezji
w~Polsce nie udostępnia do~publicznego wglądu aktów chrztów młodszych niż
100~lat oraz aktów ślubów i~zgonów młodszych niż 80~lat -- ze względu
na~prywatność danych osobowych.}. Na~podstawie grobów rodziny Gomulskich,
znajdujących się na~cmentarzu parafialnym parafii pod wezwaniem św.~Kazimierza
w~Pruszkowie, można jedynie stwierdzić, iż~urodzili się oni odpowiednio
w~1923~(patrz: ryc. \ref{fig:mgomulska_1951}) oraz w~1929~roku (patrz: ryc. 
\ref{fig:czgomulski_1946})\footnote{Drzewo genealogiczne najbliższej rodziny
Feliksa i~Małgorzaty Gomulskich zostało zaprezentowane
w~\hyperref[sec:drzewo_agumulski]{załączniku numer I}~do niniejszej książki.}.

\begin{figure}[!ht]
    \vspace*{0.5cm}
    \centering \includegraphics[width=0.83\linewidth]{
        1918_Zygmunt_Gomulski_akt_chrztu_parafia_Mińsk_Mazowiecki_wpis_196.jpg}
    \captionsetup{format=hang}
    \caption{Akt chrztu Zygmunta Gomulskiego -- par. Mińsk Mazowiecki 
    1918~rok (196/1918) 
    \cite{par_minsk2}.}
    \label{fig:zgomulski_1918}
\end{figure}

\begin{figure}[!ht]
    \vspace*{0.5cm}
    \centering \includegraphics[width=0.83\linewidth]{
        1920_Janina_Gomulska_akt_chrztu_parafia_Mińsk_Mazowiecki_wpis_347.jpg}
    \captionsetup{format=hang}
    \caption{Akt chrztu Janiny Gomulskiej -- par. Mińsk Mazowiecki 
    1920~rok (347/1920) 
    \cite{par_minsk2}.}
    \label{fig:jgomulska_1920}
\end{figure}

\begin{figure}[!ht]
    \vspace*{0.5cm}
    \centering \includegraphics[width=1.0\linewidth]{
        1917_Helena_Gomulska_akt_zgonu_parafia_Mińsk_Mazowiecki_wpis_303.jpg}
    \captionsetup{format=hang}
    \caption{Akt zgonu Heleny Gomulskiej -- par. Mińsk Mazowiecki 
    1917~rok (303/1917) 
    \cite{par_minsk2}.}
    \label{fig:hgomulska_1917_2}
\end{figure}

\begin{figure}[!ht]
    \vspace*{0.75cm}
    \centering \includegraphics[width=0.7\linewidth]{
        1941_Feliks_Gomulski_zdjęcie_nagrobka1_cmentarz_Pruszków.jpg}
    \captionsetup{format=hang}
    \caption{Zdjęcie grobu Feliksa Gomulskiego, znajdującego się na
    cmentarzu parafialnym parafii pod wezwaniem św. Kazimierza w~Pruszkowie 
    (sektor: P2, rząd: 2, numer: 59)}
    \label{fig:fgomulski_1941}
\end{figure}

Feliks i~Małgorzata Gomulscy, jak można wywnioskować na podstawie aktów 
chrztów ich dzieci, bardzo często zmieniali miejsca swojego zamieszkania. 
Były to często przeprowadzki na odległości kilkudziesięciu kilometrów, 
z~jednej strony Warszawy na drugą. Wygląda na~to, że w~latach 30. XX wieku 
przeprowadzili się ostatecznie razem z~dziećmi do Pruszkowa, gdzie żyli do 
końca swoich dni -- jak wynika z~informacji znajdujących się na ich nagrobkach, 
Feliks Gomulski\indeksosoby{Feliks}{Gomulski}{1875}{Anielina} zmarł
w~1941~roku (patrz: ryc. \ref{fig:fgomulski_1941}), natomiast jego żona
Małgorzata
z~Makowskich\indeksosoby[Makowska]{Małgorzata}{Gomulska}{1888}{Słomin} zmarła
w~1951~roku (patrz: ryc. \ref{fig:mgomulska_1951}).

\begin{figure}[!ht]
    \vspace*{0.4cm}
    \centering \includegraphics[width=0.85\linewidth]{
        1951_Małgorzata_Gomulska_Makowska_zdjęcie_nagrobka2_cmentarz_Pruszków.jpg}
    \captionsetup{format=hang}
    \caption{Zdjęcie grobu Małgorzaty, Waleriana i~Jadwigi Gomulskich, 
    znajdującego się na cmentarzu parafialnym parafii pod wezwaniem św. 
    Kazimierza w~Pruszkowie (sektor: P4, rząd: 3, numer: 21)}
    \label{fig:mgomulska_1951}
\end{figure}

\begin{figure}[!ht]
    \vspace*{0.5cm}
    \centering \includegraphics[width=0.85\linewidth]{
        1946_Czesław_Gomulski_zdjęcie_nagrobka1_cmentarz_Pruszków.jpg}
    \captionsetup{format=hang}
    \caption{Zdjęcie grobu Czesława i~Mieczysława Gomulskich, znajdującego się 
    na~cmentarzu parafialnym parafii pod wezwaniem św.~Kazimierza 
    w~Pruszkowie (sektor: B3, rząd: 4, numer: 10)}
    \label{fig:czgomulski_1946}
\end{figure}

\begin{figure}[!ht]
    \vspace*{0.5cm}
    \centering \includegraphics[width=0.8\linewidth]{
        1943_Zygmunt_Gomulski_zdjęcie_nagrobka2_cmentarz_Pruszków.jpg}
    \captionsetup{format=hang}
    \caption{Zdjęcie grobu Zygmunta Gomulskiego i~Janiny Kubiciel,
    znajdującego się na~cmentarzu parafialnym parafii pod wezwaniem
    św.~Kazimierza w~Pruszkowie (sektor: P2, rząd: 6, numer: 37)}
    \label{fig:zgomulski_1943}
\end{figure}

Dalsze losy dzieci Feliksa\indeksosoby{Feliks}{Gomulski}{1875}{Anielina}
i~Małgorzaty\indeksosoby[Makowska]{Małgorzata}{Gomulska}{1888}{Słomin}
Gomulskich nie są~niestety autorowi niniejszej książki znane. Jedyne informacje
jakie udało się pozyskać, to~te~pochodzące z~nagrobków znajdujących się
na~cmentarzu parafialnym parafii pod wezwaniem św.~Kazimierza w~Pruszkowie.
Wskazują one, iż Zygmunt Gomulski\indeksosoby{Zygmunt}{Gomulski}{1918}{Anielina}
zmarł w~1943~roku (patrz: ryc. \ref{fig:zgomulski_1943}), Czesław
Gomulski\indeksosoby{Czesław}{Gomulski}{1912}{Słomin} zmarł
w~1946~roku (patrz: ryc.\ref{fig:czgomulski_1946}), Mieczysław
Gomulski\indeksosoby{Mieczysław}{Gomulski}{1929}{Anielina} zmarł
w~1968~roku (patrz: ryc. \ref{fig:czgomulski_1946}), Jadwiga
Gomulska\indeksosoby{Jadwiga}{Gomulska}{1923}{Anielina} zmarła
w~1994~roku (patrz: ryc. \ref{fig:mgomulska_1951}), Walerian
Gomulski\indeksosoby{Walerian \flagaPL}{Gomulski}{1914}{Mińsk Mazowiecki} zmarł
w~1995~roku (patrz: ryc. \ref{fig:mgomulska_1951}), a~Janina
Gomulska\indeksosoby[Kubiciel]{Janina}{Gomulska}{1920}{Mińsk Mazowiecki}
dożyła do~2008~roku (patrz: ryc. \ref{fig:zgomulski_1943}). Co~ciekawe,
na~stronie Muzeum Powstania Warszawskiego
(\url{https://www.1944.pl}), widnieje informacja, iż~Walerian
Gomulski\indeksosoby{Walerian \flagaPL}{Gomulski}{1914}{Mińsk Mazowiecki}
był uczestnikiem Powstania Warszawskiego, gdzie walczył w~stopniu
strzelca, po~czym trafił do~niemieckiej niewoli -- numer jeniecki:~45706 --
patrz \ref{fig:wgomulski_biogram}\footnote{Biogram Powstańczy:
\url{https://www.1944.pl/powstancze-biogramy/walerian-gomulski,14633.html}}.
Historia rodziny Gomulskich w~Anielinie zakończyłaby się w~latach 30.~XX~wieku,
kiedy to~Feliks Gomulski\indeksosoby{Feliks}{Gomulski}{1875}{Anielina} wraz
z~rodziną ostatecznie wyprowadził się z~tej miejscowości (przenosząc się
prawdopodobnie do~Pruszkowa), gdyby nie fakt, iż~około 1918~roku do~Anieliny
z~Desna powrócił, wraz z~częścią swojej najbliższej rodziny, Stanisław
Gomulski\indeksosoby{Stanisław}{Gomulski}{1863}{Anielina} (ur.~1863~r.~--
zm.~1929~r.), najmłodszy syn Piotra Franciszka Gumulskiego.

\begin{figure}[!ht]
    \centering \includegraphics[width=1.0\linewidth]{
        Walerian_Gomulski_biogram_powstańczy.png}
    \captionsetup{format=hang}
    \caption{Biogram Powstańczy Waleriana Gomulskiego ze strony internetowej
    Muzeum Powstania Warszawskiego.}
    \label{fig:wgomulski_biogram}
\end{figure}

Stanisław Gomulski\indeksosoby{Stanisław}{Gomulski}{1863}{Anielina} wrócił
do~miejscowości swoich narodzin wraz z~żoną
Elżbietą\indeksosoby[Brodzińska]{Elżbieta}{Gomulska}{1857}{Grębiszew}
(ur.~1857~r. -- zm.~1928~r.), córkami
Stanisławą\indeksosoby[Gomulska]{Stanisława}{Bronikowska}{1890}{Anielina}
(ur.~1890~r. -- zm.~1976~r.)
i~Katarzyną\indeksosoby[Gomulska]{Katarzyna}{Zawadzka}{1895}{Desno}
(ur.~1895~r. -- zm.~1987~r.) oraz z~najmłodszym synem
Władysławem\indeksosoby{Władysław}{Gomulski}{1901}{Desno} (ur.~1901~r.
-- zm.~1967~r.). Pozostałe dzieci Stanisława i~Elżbiety Gomulskich, w~1918~roku,
miały już własne rodziny, toteż nie wróciły z~rodzicami do~Anieliny. Katarzyna
Gomulska w~1919~roku wyszła za mąż za Stanisława
Zawadzkiego\indeksosoby{Stanisław}{Zawadzki}{1898}{Anielina} (ur.~1898~r.~--
zm.~1982~r.) z~Anieliny. Z~ich związku zrodziło się tylko jedno dziecko --
córka Marianna\indeksosoby{Marianna}{Zawadzka}{1924}{Anielina} (ur.~1924~r. --
zm.~2015~r.). Stanisława Gomulska w~1920~roku wzięła ślub ze~Stanisławem
Bronikowskim\indeksosoby{Stanisław}{Bronikowski}{1887}{Wólka Mińska}
(ur.~1887~r. -- zm.~1955~r.) z~Wólki Mińskiej, para dochowała się trójki
dzieci: Wacławy\indeksosoby{Wacława}{Bronikowska}{1921}{Anielina}
(ur.~1921~r.~-- zm.~1986~r.),
Genowefy\indeksosoby{Genowefa}{Bronikowska}{1924}{Anielina} (ur.~1924~r. --
zm.~2015~r.) oraz
Stanisława\indeksosoby{Stanisław}{Bronikowski}{1931}{Anielina} (ur.~1931~r. --
zm.~2009~r.). Władysław Gomulski\indeksosoby{Władysław}{Gomulski}{1901}{Desno}
natomiast, w~1921~roku ożenił się z~Marianną
Paszkowską\indeksosoby[Paszkowska]{Marianna}{Gomulska}{1903}{Anielina}
(ur.~1903~r.~-- zm.~1968~r.) z~Anieliny -- wspólnie dochowali się czwórki
dzieci: Jerzego\indeksosoby{Jerzy}{Gomulski}{1922}{Anielina}
(ur.~1922~r. -- zm.~1996~r.),
Krystyny\indeksosoby[Gomulska]{Krystyna}{Krawczyńska}{1924}{Anielina}
(ur.~1924~r. -- zm.~2007~r.),
Szczepana\indeksosoby{Szczepan}{Gomulski}{1927}{Anielina} (ur.~1927~r. --
zm.~2014~r.) oraz
Benedykta\indeksosoby{Benedykt}{Gomulski}{1938}{Anielina} (ur.~1938~r.~--
zm.~2021~r.) -- wszystkie z~nich urodziły się w~Anielinie, a~ich potomkowie
mieszkają tam prawdopodobnie do dnia dzisiejszego\footnote{Oczywiście już nie
w~Anielinie, bo~od~1~stycznia 1986~roku Anielina stała się częścią Mińska
Mazowieckiego i~jedyny widoczny ślad po~tej miejscowości stanowi przystanek
kolejowy Mińsk Mazowiecki Anielina.}.

% Przednia okładka podrozdziału
\newpage
\includepdf{Wolka_Minska_mapa_fin.png}

\phantomsection
\addcontentsline{toc}{section}{Wólka Mińska: 1861~r. -- \biezacyrok}

\lettrine[lines=5, lhang=0.1]{W}{ólka} Mińska to~wieś szlachecka założona
w~połowie XVI~wieku w~odległości około trzech kilometrów na~północ od~Mińska
Mazowieckiego, na~terenie ziemi czerskiej województwa mazowieckiego.
Od~początku swojego istnienia wieś ta~należała do~parafii pod wezwaniem
Narodzenia Najświętszej Maryi Panny w~Mińsku Mazowieckim. Na~początku
XIX~wieku Wólka Mińska znajdowała się na~terenie zaboru austriackiego,
następnie w~1809~roku została włączona do~Księstwa Warszawskiego,
a~w~1815~roku do~Królestwa Polskiego. Słownik geograficzny Królestwa Polskiego
podaje, iż~w~1827~roku w~Wólce Mińskiej znajdowały się trzy domy, w~których
mieszkało 25~osób. Pierwszy powszechny spis ludności po~odzyskaniu przez
Polskę niepodległości wykazał, iż~w~1921~roku w~Wólce Mińskiej było 10~domów,
w~których mieszkały 84~osoby. Nieopodal Wólki Mińskiej swoje źródła ma~rzeka
Długa\footnote{Rzeka Długa przepływa również niedaleko wsi Desno, co~więcej
odegrała ona prawdopodobnie istotną rolę w~zdefiniowaniu nazwy tej
miejscowości, o~czym szerzej opowiem w~dalszej części niniejszej książki.}.

\begin{figure}[!ht]
    \vspace*{0.5cm}
    \centering \includegraphics[width=0.75\linewidth]{
        Wólka Mińska - Mazowsze w drugiej połowie XVI wieku.png}
    \captionsetup{format=hang}
    \caption{Mapa miejscowości znajdujących się w~okolicy Mińska Mazowieckiego 
    w~połowie XVI~wieku opracowana w~1972~roku przez Instytut Historii 
    Polskiej Akademii Nauk \cite{palucki}.}
    \label{fig:wolka_minska_xvi}
\end{figure}

Pierwszym znanym członkiem rodziny Gomulskich, który zamieszkał w~Wólce
Mińskiej był Jan Gomulski\indeksosoby{Jan}{Gomulski}{1836}{Zamienie}
(ur.~1836~r. -- zm.~1901~r.), najstarszy syn Piotra
Franciszka\indeksosoby{Piotr Franciszek}{Gumulski}{1815}{Rudzienko}
i~Marianny Gumulskich\indeksosoby[Chłopik]{Marianna}{Gumulska}{1819}{Zamienie}.
Dnia 31~stycznia~1858~roku wziął on~ślub z~Agnieszką
Piotrkowicz\indeksosoby[Piotrkowicz]{Agnieszka}{Gomulska}{1840}{Budy
Barcząckie} (ur.~1840~r.~-- zm.~1881~r.) -- patrz: ryc.
\ref{fig:jgomulski_1858}, która urodziła się 26~grudnia~1840~roku w~Budach
Barcząckich (patrz: ryc. \ref{fig:agumulska_1840}). Rodzicami Agnieszki byli
Tomasz\indeksosoby{Tomasz}{Piotrkowicz}{1804}{Dębe Wielkie} (ur.~1804~r.~--
zm.~1874~r.) oraz
Marianna\indeksosoby[Lisiewska]{Marianna}{Piotrkowicz}{1811}{Dębe Wielkie} 
(ur.~1811~r.~-- zm.~1874~r.)z~domu Lisiewska, pochodzili oni z~Dębego
(obecnie Dębe Wielkie)\footnote{Co ciekawe pradziadek Agnieszki Gomulskiej
z~domu Piotrkowicz, Maciej
Okrzeja\indeksosoby{Maciej}{Okrzeja}{1752}{Dębe Wielkie}, jest również
praprapradziadkiem Prezesa Rady Ministrów Jana Ferdynanda
Olszewskiego\indeksosoby{Jan Ferdynand \flagaPL}{Olszewski}{1930}{Warszawa}
(ur.~1930~r. -- zm.~2019~r.) oraz prapradziadkiem Stefana Aleksandra
Okrzei\indeksosoby{Stefan Aleksander \flagaPL}{Okrzeja}{1886}{Dębe Wielkie}
(ur.~1886~r. -- zm.~1905~r.), działacza niepodległościowego i~socjalistycznego,
co~zostało szerzej opisane w~\hyperref[sec:piotrkowicze]{załączniku numer~V}
do~niniejszej książki.}.

\begin{figure}[!ht]
    \vspace*{0.5cm}
    \centering \includegraphics[width=0.95\linewidth]{
        1840_Agnieszka_Gumulska_Piotrkowicz_akt_chrztu_parafia_Mińsk_Mazowiecki_wpis_314.jpg}
    \captionsetup{format=hang}
    \caption{Akt chrztu Agnieszki Piotrkowicz -- par. Mińsk Mazowiecki 
    1840~rok (314/1840) 
    \cite{par_minsk2}.}
    \label{fig:agumulska_1840}
\end{figure}

Przez pierwsze lata po ślubie Jan\indeksosoby{Jan}{Gomulski}{1836}{Zamienie}
i~Agnieszka\indeksosoby[Piotrkowicz]{Agnieszka}{Gomulska}{1840}{Budy
Barcząckie} Gomulscy zamieszkiwali jeszcze w~rodzinnej wsi Jana -- w~Anielinie,
gdzie 22~stycznia 1859~roku urodził się ich pierwszy syn
Józef\indeksosoby{Józef}{Gomulski}{1859}{Anielina} (ur.~1859~r.~-- zm.~1929~r.)
-- patrz: ryc. \ref{fig:jgomulski_1859}). W~akcie chrztu Józefa, jego ojciec
Jan określony został jako wyrobnik. W~okolicach 1861~roku Jan i~Agnieszka
Gomulscy wraz z~synem Józefem przeprowadzili się do~Wólki Mińskiej, gdzie
22~lutego 1861~roku urodził się ich drugi syn --
Stanisław\indeksosoby{Stanisław}{Gomulski}{1861}{Wólka Mińska} (ur.~1861~r.~--
zm.~1910~r.). Rodzicami chrzestnymi Stanisława zostali Kacper Bronikowski oraz
Marianna Frelak, a~jego ojciec w~akcie chrztu został określony jako \enquote{...
gospodarz, ojciec dziecięcia lat dwadzieścia pięć mający w~Wólce
zamieszkały...} (patrz: ryc. \ref{fig:sgomulski_1861}). Stanisław Gomulski był
ostatnim z~trzech członków rodziny Gomulskich, którzy byli współzałożycielami
wsi Desna (a~zarazem udziałowcami spółki \enquote{Desna}) -- razem ze~swoimi
stryjkami (a~jednocześnie rówieśnikami)
Michałem\indeksosoby{Michał}{Gomulski}{1859}{Anielina}
i~Stanisławem\indeksosoby{Stanisław}{Gomulski}{1863}{Anielina} (stryj Michał
był od niego o~dwa lata starszy, a~jego stryjk-imiennik o~dwa lata młodszy)
przeprowdził się do~wsi Desna w~okoliach 1896~roku. Bezpośredni potomkowie
Stanisława Gomulskiego mieszkają we~wsi Desno do~dnia dzisiejszego, nosząc
nazwisko po~swoim
przodku\footnote{Autor\indeksosoby{Mateusz}{Gomulski}{1990}{Warszawa}
niniejszej książki jest praprawnukiem Stanisława Gomulskiego, mieszka
we~wsi Desno, a~ponadto urodził się dokładnie 129~lat po~nim – 22~lutego
1990~roku.}.

\begin{figure}[!ht]
    \vspace*{0.5cm}
    \centering \includegraphics[width=0.92\linewidth]{
        1858_Jan_Gomulski_Agnieszka_Piotrkowicz_akt_ślubu_parafia_Mińsk_Mazowiecki_wpis_14.jpg}
    \captionsetup{format=hang}
    \caption{Akt ślubu Jana Gomulskiego i~Agnieszki Piotrkowicz -- par. 
    Mińsk Mazowiecki 1858~rok (14/1858) \cite{par_minsk2}.}
    \label{fig:jgomulski_1858}
\end{figure}

\begin{figure}[!ht]
    \vspace*{0.5cm}
    \centering \includegraphics[width=1.0\linewidth]{
        1859_Józef_Gomulski_akt_chrztu_parafia_Mińsk_Mazowiecki_wpis_25.jpg}
    \captionsetup{format=hang}
    \caption{Akt chrztu Józefa Gomulskiego -- par. Mińsk Mazowiecki 
    1859~rok (25/1859) 
    \cite{par_minsk2}.}
    \label{fig:jgomulski_1859}
\end{figure}

\begin{figure}[!ht]
    \vspace*{0.5cm}
    \centering \includegraphics[width=1.0\linewidth]{
        1861_Stanisław_Gomulski_akt_chrztu_parafia_Mińsk_Mazowiecki_wpis_56.jpg}
    \captionsetup{format=hang}
    \caption{Akt chrztu Stanisława Gomulskiego -- par. Mińsk Mazowiecki 
    1861~rok (56/1861) 
    \cite{par_minsk2}.}
    \label{fig:sgomulski_1861}
\end{figure}

Jan\indeksosoby{Jan}{Gomulski}{1836}{Zamienie}
i~Agnieszka\indeksosoby[Piotrkowicz]{Agnieszka}{Gomulska}{1840}{Budy
Barcząckie} Gomulscy przez 23~lata małżeństwa doczekali się łącznie
dwanaściorga dzieci (wszystkie ich dzieci, poza Józefem, urodziły się w~Wólce
Mińskiej):

\begin{itemize}
    \item Józef Gomulski\indeksosoby{Józef}{Gomulski}{1859}{Anielina}
    (ur. 1859~r. -- zm. 1929~r.) -- patrz: ryc. \ref{fig:jgomulski_1859},
    \item Stanisław Gomulski\indeksosoby{Stanisław}{Gomulski}{1861}{Wólka Mińska}
    (ur. 1861~r. -- zm. 1910~r.) -- patrz: ryc. \ref{fig:sgomulski_1861},
    \item Jan Gomulski\indeksosoby{Jan}{Gomulski}{1863}{Wólka Mińska}
    (ur. 1863~r. -- zm. 1932~r.) -- patrz: ryc. \ref{fig:jgomulski_1863},
    \item Julianna Gomulska\indeksosoby{Julianna}{Gomulska}{1865}{Wólka Mińska}
    (ur. 1865~r. -- zm. 1868~r.) -- patrz: ryc. \ref{fig:jgomulska_1865},
    \item Marianna Gomulska\indeksosoby{Marianna}{Gomulska}{1867}{Wólka Mińska}
    (ur. 1867~r. -- zm. 1868~r.) -- patrz: ryc. \ref{fig:mgomulska_1867},
    \item Józefa Gomulska\indeksosoby{Józefa}{Gomulska}{1869}{Wólka Mińska}
    (ur. 1869~r. -- zm. 1872~r.) -- patrz: ryc. \ref{fig:jgomulska_1869},
    \item Anna Gomulska\indeksosoby[Gomulska]{Anna}{Kwiatkowska}{1872}{Wólka Mińska}
    (ur. 1872~r. -- zm. ?) -- patrz: ryc. \ref{fig:agomulska_1872},
    \item Antonina Paulina Gomulska\indeksosoby{Antonina Paulina}{Gomulska}{1874}{Wólka Mińska}
    (ur. 1874~r. -- zm. 1877~r.) -- patrz: ryc. \ref{fig:apgomulska_1874},
    \item Zofia Gomulska\indeksosoby[Gomulska]{Zofia}{Wojdyga}{1876}{Wólka Mińska}
    (ur. 1876~r. -- zm. 1948~r.) -- patrz: ryc. \ref{fig:zgomulska_1876},
    \item Franciszek Gomulski\indeksosoby{Franciszek}{Gomulski}{1879}{Wólka Mińska}
    (ur. 1879~r. -- zm. 1886~r.) -- patrz: ryc. \ref{fig:fgomulski_1879},
    \item Antoni Gomulski\indeksosoby{Antoni}{Gomulski}{1879}{Wólka Mińska}
    (ur. 1879~r. -- zm. ?) -- patrz: ryc. \ref{fig:agomulski_1879},
    \item Marianna Gomulska\indeksosoby{Marianna}{Gomulska}{1881}{Wólka Mińska}
    (ur. 1881~r. -- zm. 1881~r.) -- patrz: ryc. \ref{fig:mgomulska_1881}.
  \end{itemize}

Tylko szóstka dzieci małżeństwa Gomulskich dożyła wieku dorosłego: Józef,
Stanisław, Jan, Anna, Zofia oraz Antoni.

\begin{figure}[!ht]
    \vspace*{0.5cm}
    \centering \includegraphics[width=0.9\linewidth]{
        1863_Jan_Gomulski_akt_chrztu_parafia_Mińsk_Mazowiecki_wpis_140.jpg}
    \captionsetup{format=hang}
    \caption{Akt chrztu Jana Gomulskiego -- par. Mińsk Mazowiecki 1863~rok (140/1863) 
    \cite{par_minsk2}.}
    \label{fig:jgomulski_1863}
\end{figure}

\begin{figure}[!ht]
    \vspace*{0.5cm}
    \centering \includegraphics[width=0.95\linewidth]{
        1865_Julianna_Gomulska_akt_chrztu_parafia_Mińsk_Mazowiecki_wpis_81.jpg}
    \captionsetup{format=hang}
    \caption{Akt chrztu Julianny Gomulskiej -- par. Mińsk Mazowiecki 1865~rok
    (81/1865) \cite{par_minsk2}.}
    \label{fig:jgomulska_1865}
\end{figure}

\begin{figure}[!ht]
    \vspace*{0.5cm}
    \centering \includegraphics[width=0.95\linewidth]{
        1867_Marianna_Gomulska_akt_chrztu_parafia_Mińsk_Mazowiecki_wpis_60.jpg}
    \captionsetup{format=hang}
    \caption{Akt chrztu Marianny Gomulskiej -- par. Mińsk Mazowiecki 1867~rok
    (60/1867) \cite{par_minsk2}.}
    \label{fig:mgomulska_1867}
\end{figure}

\begin{figure}[!ht]
    \vspace*{0.5cm}
    \centering \includegraphics[width=0.95\linewidth]{
        1869_Józefa_Gomulska_akt_chrztu_parafia_Mińsk_Mazowiecki_wpis_91.jpg}
    \captionsetup{format=hang}
    \caption{Akt chrztu Józefy Gomulskiej -- par. Mińsk Mazowiecki 1869~rok
    (91/1869) \cite{par_minsk2}.}
    \label{fig:jgomulska_1869}
\end{figure}

\begin{figure}[!ht]
    \vspace*{0.5cm}
    \centering \includegraphics[width=0.77\linewidth]{
        1872_Anna_Gomulska_Kwiatkowska_akt_chrztu_parafia_Mińsk_Mazowiecki_wpis_98.jpg}
    \captionsetup{format=hang}
    \caption{Akt chrztu Anny Gomulskiej -- par. Mińsk Mazowiecki 1872~rok
    (98/1872) \cite{par_minsk2}.}
    \label{fig:agomulska_1872}
\end{figure}

\begin{figure}[!ht]
    \vspace*{0.5cm}
    \centering \includegraphics[width=0.77\linewidth]{
        1874_Antonina_Paulina_Gomulska_akt_chrztu_parafia_Mińsk_Mazowiecki_wpis_122.jpg}
    \captionsetup{format=hang}
    \caption{Akt chrztu Antoniny Pauliny Gomulskiej -- par. Mińsk Mazowiecki
    1874~rok (122/1874) 
    \cite{par_minsk2}.}
    \label{fig:apgomulska_1874}
\end{figure}

\begin{figure}[!ht]
    \vspace*{0.5cm}
    \centering \includegraphics[width=0.82\linewidth]{
        1876_Zofia_Gomulska_Wojdyga_akt_chrztu_parafia_Mińsk_Mazowiecki_wpis_158.jpg}
    \captionsetup{format=hang}
    \caption{Akt chrztu Zofii Gomulskiej -- par. Mińsk Mazowiecki 1876~rok
    (158/1876) \cite{par_minsk2}.}
    \label{fig:zgomulska_1876}
\end{figure}

\begin{figure}[!ht]
    \vspace*{0.5cm}
    \centering \includegraphics[width=0.82\linewidth]{
        1879_Franciszek_Gomulski_akt_chrztu_parafia_Mińsk_Mazowiecki_wpis_15.jpg}
    \captionsetup{format=hang}
    \caption{Akt chrztu Franciszka Gomulskiego -- par. Mińsk Mazowiecki
    1879~rok (15/1879) 
    \cite{par_minsk2}.}
    \label{fig:fgomulski_1879}
\end{figure}

\begin{figure}[!ht]
    \vspace*{0.5cm}
    \centering \includegraphics[width=0.82\linewidth]{
        1879_Antoni_Gomulski_akt_chrztu_parafia_Mińsk_Mazowiecki_wpis_16.jpg}
    \captionsetup{format=hang}
    \caption{Akt chrztu Antoniego Gomulskiego -- par. Mińsk Mazowiecki
    1879~rok (16/1879) 
    \cite{par_minsk2}.}
    \label{fig:agomulski_1879}
\end{figure}

\begin{figure}[!ht]
    \vspace*{0.5cm}
    \centering \includegraphics[width=0.82\linewidth]{
        1881_Marianna_Gomulska_akt_chrztu_parafia_Mińsk_Mazowiecki_wpis_124.jpg}
    \captionsetup{format=hang}
    \caption{Akt chrztu Marianny Gomulskiej -- par. Mińsk Mazowiecki
    1881~rok (124/1881) 
    \cite{par_minsk2}.}
    \label{fig:mgomulska_1881}
\end{figure}

Agnieszka Gomulska\indeksosoby[Piotrkowicz]{Agnieszka}{Gomulska}{1840}{Budy
Barcząckie} z~domu Piotrkowicz zmarła w~Wólce Mińskiej 27~kwietnia
1881~roku, w~wieku 40~lat (patrz: ryc. \ref{fig:agomulska_1881}). Przyczyną
jej śmierci były prawdopodobnie powikłania po~porodowe, gdyż kilka dni
wcześniej, 19~kwietnia 1881~roku, urodziła ona swoją najmłodszą córkę Mariannę,
która również zmarła tego samego dnia co~jej matka.

\begin{figure}[!ht]
    \vspace*{0.5cm}
    \centering \includegraphics[width=0.75\linewidth]{
        1881_Agnieszka_Gumulska_Piotrkowicz_akt_zgonu_parafia_Mińsk_Mazowiecki_wpis_105.jpg}
    \captionsetup{format=hang}
    \caption{Akt zgonu Agnieszki Gomulskiej (z~domu Piotrkowicz) -- par. Mińsk
    Mazowiecki 1881~rok (105/1881) 
    \cite{par_minsk2}.}
    \label{fig:agomulska_1881}
\end{figure}

Jan Gomulski\indeksosoby{Jan}{Gomulski}{1836}{Zamienie} niecałe dwa miesiące
po~śmierci swojej żony
Agnieszki\indeksosoby[Piotrkowicz]{Agnieszka}{Gomulska}{1840}{Budy
Barcząckie}, 16~czerwca 1881~roku, ożenił się ponownie --
z~26~lat młodszą, niespełna dwudziestoletnią, Wiktorią
Świętochowską\indeksosoby[Świętochowska]{Wiktoria}{Gomulska}{1862}{Dłużka}
(ur. 1862~r. -- zm. 1937~r.) z~pobliskiej wsi Dłużka (patrz: ryc.
\ref{fig:jgomulski_1881})\footnote{Co~ciekawe trzeci syn
Jana\indeksosoby{Jan}{Gomulski}{1836}{Zamienie} z~pierwszego małżeństwa,
również Jan\indeksosoby{Jan}{Gomulski}{1863}{Wólka Mińska}, zaledwie pieć
miesięcy później ożenił się z~młodszą siostrą swojej macochy -- Marianną
Świętochowską\indeksosoby[Świętochowska]{Marianna}{Gomulska}{1864}{Dłużka}
(ur. 1864~r. -- zm. 1912~r.).}. Jan i~Wiktoria Gomulscy doczekali się razem
jedenaściorga dzieci:

\begin{itemize}
    \item Marianna Gomulska\indeksosoby[Gomulska]{Marianna}{Królak}{1882}{Wólka
    Mińska} (ur. 1882~r. -- zm. 1953~r.) -- patrz: ryc. \ref{fig:mgomulska_1882},
    \item Józef Gomulski\indeksosoby{Józef}{Gomulski}{1884}{Mińsk Mazowiecki}
    (ur. 1884~r. -- zm. 1929~r.) -- patrz: ryc. \ref{fig:jgomulski_1884},
    \item Ludwik Gomulski\indeksosoby{Ludwik}{Gomulski}{1886}{Mińsk Mazowiecki}
    (ur. 1886~r. -- zm. ?) -- patrz: ryc. \ref{fig:lgomulski_1886},
    \item Franciszek Stefan Gomulski\indeksosoby{Franciszek
    Stefan}{Gomulski}{1888}{Mińsk Mazowiecki} (ur. 1888~r. -- zm. 1889~r.) --
    patrz: ryc. \ref{fig:fsgomulski_1889},
    \item Bezimienna Gomulska\indeksosoby{Bezimienna}{Gomulska}{1891}{Mińsk
    Mazowiecki} (ur. 1891~r. -- zm. 1891~r.) -- patrz: ryc.
    \ref{fig:bgomulska_1891},
    \item Władysław Gomulski\indeksosoby{Władysław}{Gomulski}{1892}{Mińsk
    Mazowiecki} (ur. 1892~r. -- zm. 1893~r.) -- patrz: ryc. 
    \ref{fig:wgomulski_1892},
    \item Eleonora Gomulska\indeksosoby[Gomulska][Klimek]{Eleonora}{Wilk}{1894}{Wólka
    Mińska} (ur. 1894~r. -- zm. 1957~r.) -- patrz: ryc. \ref{fig:egomulska_1894},
    \item Mieczysław Gomulski\indeksosoby{Mieczysław}{Gomulski}{1896}{Wólka
    Mińska}(ur. 1896~r. -- zm. 1967~r.) -- patrz: ryc. \ref{fig:mgomulski_1896},
    \item Wanda Gomulska\indeksosoby{Wanda}{Gomulska}{1897}{Wólka
    Mińska} (ur. 1897~r. -- zm. 1898~r.) -- patrz: ryc. \ref{fig:wgomulska_1898},
    \item Bronisław Gomulski\indeksosoby{Bronisław}{Gomulski}{1899}{Wólka
    Mińska} (ur. 1899~r. -- zm. 1905~r.) -- patrz: ryc. \ref{fig:bgomulski_1899},
    \item Wanda Gomulska\indeksosoby[Gomulska]{Wanda}{Majszyk}{1901}{Wólka
    Mińska} (ur. 1901~r. -- zm. 1978~r.) -- patrz: ryc. \ref{fig:wgomulska_1901}.
  \end{itemize}

  \begin{figure}[!ht]
    \vspace*{0.5cm}
    \centering \includegraphics[width=0.85\linewidth]{
        1881_Jan_Gumulski_Wiktoria_Świętochowska_akt_ślubu_parafia_Mińsk_Mazowiecki_wpis_33.jpg}
    \captionsetup{format=hang}
    \caption{Akt ślubu Jana Gomulskiego i~Wiktorii Świętochowskiej -- par. 
    Mińsk Mazowiecki 1881~rok (33/1881) \cite{par_minsk2}.}
    \label{fig:jgomulski_1881}
\end{figure}

\begin{figure}[!ht]
    \vspace*{0.5cm}
    \centering \includegraphics[width=0.77\linewidth]{
        1882_Marianna_Gomulska_Królak_akt_chrztu_parafia_Mińsk_Mazowiecki_wpis_125.jpg}
    \captionsetup{format=hang}
    \caption{Akt chrztu Marianny Gomulskiej -- par. Mińsk Mazowiecki
    1882~rok (125/1882)
    \cite{par_minsk2}.}
    \label{fig:mgomulska_1882}
\end{figure}

\begin{figure}[!ht]
    \vspace*{0.5cm}
    \centering \includegraphics[width=0.77\linewidth]{
        1884_Józef_Gomulski_akt_chrztu_parafia_Mińsk_Mazowiecki_wpis_382.jpg}
    \captionsetup{format=hang}
    \caption{Akt chrztu Józefa Gomulskiego -- par. Mińsk Mazowiecki
    1884~rok (382/1884) 
    \cite{par_minsk2}.}
    \label{fig:jgomulski_1884}
\end{figure}

\begin{figure}[!ht]
    \vspace*{0.5cm}
    \centering \includegraphics[width=0.8\linewidth]{
        1886_Ludwik_Gomulski_akt_chrztu_parafia_Mińsk_Mazowiecki_wpis_386.jpg}
    \captionsetup{format=hang}
    \caption{Akt chrztu Ludwika Gomulskiego -- par. Mińsk Mazowiecki
    1886~rok (386/1886) 
    \cite{par_minsk2}.}
    \label{fig:lgomulski_1886}
\end{figure}

\begin{figure}[!ht]
    \vspace*{0.5cm}
    \centering \includegraphics[width=0.8\linewidth]{
        1889_Franciszek_Stefan_Gomulski_akt_chrztu_parafia_Mińsk_Mazowiecki_wpis_8.jpg}
    \captionsetup{format=hang}
    \caption{Akt chrztu Franciszka Stefana Gomulskiego -- par. Mińsk Mazowiecki
    1889~rok (8/1889) 
    \cite{par_minsk2}.}
    \label{fig:fsgomulski_1889}
\end{figure}

\begin{figure}[!ht]
    \vspace*{0.5cm}
    \centering \includegraphics[width=0.8\linewidth]{
        1891_Bezimienna_Gomulska_akt_chrztu_akt_zgonu_parafia_Mińsk_Mazowiecki_wpis_124.jpg}
    \captionsetup{format=hang}
    \caption{Akt chrztu Bezimiennej Gomulskiej -- par. Mińsk Mazowiecki
    1891~rok (124/1891) 
    \cite{par_minsk2}.}
    \label{fig:bgomulska_1891}
\end{figure}

\begin{figure}[!ht]
    \vspace*{0.5cm}
    \centering \includegraphics[width=0.8\linewidth]{
        1892_Władysław_Gomulski_akt_chrztu_parafia_Mińsk_Mazowiecki_wpis_218.jpg}
    \captionsetup{format=hang}
    \caption{Akt chrztu Władysława Gomulskiego -- par. Mińsk Mazowiecki
    1892~rok (218/1892) 
    \cite{par_minsk2}.}
    \label{fig:wgomulski_1892}
\end{figure}

\begin{figure}[!ht]
    \vspace*{0.5cm}
    \centering \includegraphics[width=0.77\linewidth]{
        1894_Eleonora_Gomulska_Wilk_akt_chrztu_parafia_Mińsk_Mazowiecki_wpis_460.jpg}
    \captionsetup{format=hang}
    \caption{Akt chrztu Eleonory Gomulskiej -- par. Mińsk Mazowiecki
    1894~rok (460/1894)
    \cite{par_minsk2}.}
    \label{fig:egomulska_1894}
\end{figure}

\begin{figure}[!ht]
    \vspace*{0.5cm}
    \centering \includegraphics[width=0.77\linewidth]{
        1896_Mieczysław_Gomulski_akt_chrztu_parafia_Mińsk_Mazowiecki_wpis_433.jpg}
    \captionsetup{format=hang}
    \caption{Akt chrztu Mieczysława Gomulskiego -- par. Mińsk Mazowiecki
    1896~rok (433/1896)
    \cite{par_minsk2}.}
    \label{fig:mgomulski_1896}
\end{figure}

\begin{figure}[!ht]
    \vspace*{0.5cm}
    \centering \includegraphics[width=0.75\linewidth]{
        1898_Wanda_Gomulska_akt_chrztu_parafia_Mińsk_Mazowiecki_wpis_28.jpg}
    \captionsetup{format=hang}
    \caption{Akt chrztu Wandy Gomulskiej -- par. Mińsk Mazowiecki
    1898~rok (28/1898)
    \cite{par_minsk2}.}
    \label{fig:wgomulska_1898}
\end{figure}

\begin{figure}[!ht]
    \vspace*{0.5cm}
    \centering \includegraphics[width=0.75\linewidth]{
        1899_Bronisław_Gomulski_akt_chrztu_parafia_Mińsk_Mazowiecki_wpis_135.jpg}
    \captionsetup{format=hang}
    \caption{Akt chrztu Bronisława Gomulskiego -- par. Mińsk Mazowiecki
    1899~rok (135/1899)
    \cite{par_minsk2}.}
    \label{fig:bgomulski_1899}
\end{figure}

\begin{figure}[!ht]
    \vspace*{0.5cm}
    \centering \includegraphics[width=0.95\linewidth]{
        1901_Wanda_Gomulska_akt_chrztu_parafia_Mińsk_Mazowiecki_wpis_370.jpg}
    \captionsetup{format=hang}
    \caption{Akt chrztu Wandy Gomulskiej -- par. Mińsk Mazowiecki
    1901~rok (370/1901)
    \cite{par_minsk2}.}
    \label{fig:wgomulska_1901}
\end{figure}

Jan\indeksosoby{Jan}{Gomulski}{1836}{Zamienie}
i~Wiktoria\indeksosoby[Świętochowska]{Wiktoria}{Gomulska}{1862}{Dłużka}
Gomulscy po~ślubie zamieszkali w~Wólce Mińskiej, gdzie w~1882~roku urodziła
się ich pierwsza córka Marianna. Od~1884~roku, przez około 10~lat mieszkali
w~Mińsku Mazowieckim, na~to~przynajmniej wskazują akty urodzenia ich kolejnych
dzieci: Józefa\indeksosoby{Józef}{Gomulski}{1884}{Mińsk Mazowiecki} (1884~r.),
Ludwika\indeksosoby{Ludwik}{Gomulski}{1886}{Mińsk Mazowiecki} (1886~r),
Franciszka Stefana\indeksosoby{Franciszek Stefan}{Gomulski}{1888}{Mińsk
Mazowiecki} (1889~r), Bezimiennej\indeksosoby{Bezimienna}{Gomulska}{1891}{Mińsk
Mazowiecki} (1891~r) oraz Władysława\indeksosoby{Władysław}{Gomulski}{1892}{Mińsk
Mazowiecki} (1892~r). Od 1894~roku dzieci Jana i~Wiktorii Gomulskich znowu
zaczynają przychodzić na świat w~Wólce Mińskiej, wygląda więc na~to,
że~powrócili oni w~tym czasie do~swojej dawnej wsi. W~1894~roku na~świat
w~Wólce Mińskiej przyszła
Eleonora\indeksosoby[Gomulska][Klimek]{Eleonora}{Wilk}{1894}{Wólka Mińska},
w~1896~roku
Mieczysław\indeksosoby{Mieczysław}{Gomulski}{1896}{Wólka Mińska}, w~1898~roku
Wanda\indeksosoby{Wanda}{Gomulska}{1897}{Wólka Mińska}, w~1899~roku
Bronisław\indeksosoby{Bronisław}{Gomulski}{1899}{Wólka Mińska} a~w~1901~roku
Wanda\indeksosoby[Gomulska]{Wanda}{Majszyk}{1901}{Wólka Mińska} (ich
poprzednia córka Wanda zmarła niedługo po~urodzeniu). Jan
Gomulski\indeksosoby{Jan}{Gomulski}{1836}{Zamienie} zmarł 6~marca 1901~roku
w~Wólce Mińskiej, w~wieku 65~lat (patrz: ryc. \ref{fig:jgomulski_1901}).
Dochował się on~łącznie dwudziestu trojga dzieci z~dwóch małżeństw,
z~Agnieszką miał dwanaścioro dzieci, a~z~Wiktorią jedenaścioro\footnote{Drzewo
genealogiczne najbliższej rodziny Jana Gomulskiego oraz jego dwóch żon zostało
zaprezentowane w~\hyperref[sec:drzewo_jgomulski]{załączniku numer I}~do
niniejszej książki.}. Pierwsze dziecko Jana, jego najstarszy syn
Józef\indeksosoby{Józef}{Gomulski}{1859}{Anielina}, urodził się gdy Jan miał
23~lata, a~ostanie jego dziecko, córka
Wanda\indeksosoby[Gomulska]{Wanda}{Majszyk}{1901}{Wólka Mińska}, urodziła się
pięć miesięcy po~jego śmierci -- 21~sierpnia 1901~roku. Wiktoria
Gomulska\indeksosoby[Świętochowska]{Wiktoria}{Gomulska}{1862}{Dłużka}
zmarła w~Wólce Mińskiej 36~lat po~swoimi mężu, 3~grudnia 1937~roku, w~wieku
75~lat (patrz: ryc. \ref{fig:wgomulska_1937}).

\begin{figure}[!ht]
    \vspace*{0.5cm}
    \centering \includegraphics[width=0.77\linewidth]{
        1901_Jan_Gomulski_akt_zgonu_parafia_Mińsk_Mazowiecki_wpis_44.jpg}
    \captionsetup{format=hang}
    \caption{Akt zgonu Jana Gomulskiego -- par. Mińsk Mazowiecki 1901~rok (44/1901) \cite{par_minsk2}.}
    \label{fig:jgomulski_1901}
\end{figure}

\begin{figure}[!ht]
    \vspace*{0.5cm}
    \centering \includegraphics[width=0.8\linewidth]{
        1937_Wiktoria_Gumulska_Świętochowska_akt_zgonu_parafia_Mińsk_Mazowiecki_wpis_344.jpg}
    \captionsetup{format=hang}
    \caption{Akt zgonu Wiktorii Gomulskiej z d. Świętochowskiej -- par. Mińsk
    Mazowiecki 1937~rok (344/1937) \cite{par_minsk2}.}
    \label{fig:wgomulska_1937}
\end{figure}

\begin{figure}[!ht]
    \vspace*{0.4cm}
    \centering \includegraphics[width=0.77\linewidth]{
        1967_Władysław_Gomulski_zdjęcie_nagrobka1_cmentarz_Mińsk_Mazowiecki.jpg}
    \captionsetup{format=hang}
    \caption{Zdjęcie grobu Marianny i~Władysława Gomulskich, znajdującego się
    na~cmentarzu parafialnym parafii pod wezwaniem Narodzenia Najświętszej
    Maryi Panny w~Mińsku Mazowieckim (sektor: X, rząd: 1, numer: 44)}
    \label{fig:wgomulski_1967}
\end{figure}

\begin{figure}[!ht]
    \vspace*{0.4cm}
    \centering \includegraphics[width=0.77\linewidth]{
        1932_Jan_Gomulski_zdjęcie_nagrobka1_cmentarz_Mińsk_Mazowiecki.jpg}
    \captionsetup{format=hang}
    \caption{Zdjęcie grobu Jana Gomulskiego, znajdującego się na cmentarzu
    parafialnym parafii podmwezwaniem Narodzenia Najświętszej Maryi Panny
    w~Mińsku Mazowieckim (sektor: XIII, rząd: 15, numer: 18)}
    \label{fig:jgomulski_1932}
\end{figure}

\begin{figure}[!ht]
    \vspace*{0.4cm}
    \centering \includegraphics[width=0.77\linewidth]{
        1996_Jerzy_Gomulski_zdjęcie_nagrobka1_cmentarz_Mińsk_Mazowiecki.jpg}
    \captionsetup{format=hang}
    \caption{Zdjęcie grobu Zofii i~Jerzego Gomulskich, znajdującego się na
    cmentarzu parafialnym parafii pod wezwaniem Narodzenia Najświętszej Maryi
    Panny w~Mińsku Mazowieckim (sektor: IX, rząd: A, numer: 22)}
    \label{fig:jgomulski_1996}
\end{figure}

\begin{figure}[!ht]
    \vspace*{0.4cm}
    \centering \includegraphics[width=0.77\linewidth]{
        1957_Mieczysław_Gomulski_zdjęcie_nagrobka1_cmentarz_Mińsk_Mazowiecki.jpg}
    \captionsetup{format=hang}
    \caption{Zdjęcie grobu Marianny i~Mieczysława Gomulskich, znajdującego się
    na~cmentarzu parafialnym parafii pod wezwaniem Narodzenia Najświętszej
    Maryi Panny w~Mińsku Mazowieckim (sektor: IX, rząd: 9, numer: 14)}
    \label{fig:mgomulski_1957}
\end{figure}

\begin{figure}[!ht]
    \vspace*{0.4cm}
    \centering \includegraphics[width=0.77\linewidth]{
        1957_Eleonora_Gomulska_Klimek_Wilk_zdjęcie_nagrobka1_cmentarz_Mińsk_Mazowiecki.jpg}
    \captionsetup{format=hang}
    \caption{Zdjęcie grobu Wiktorii Gomulskiej z~d.~Świętochowskiej, Eleonory
    Klimek-Wilk z~d.~Gomulskiej oraz Stanisława Klimka, znajdującego się na
    cmentarzu parafialnym parafii pod wezwaniem Narodzenia Najświętszej Maryi
    Panny w~Mińsku Mazowieckim (sektor: XVIII, rząd: 5, numer: 32)}
    \label{fig:egomulski_1957}
\end{figure}

Potomkowie Piotra Franciszka
Gumluskiego\indeksosoby{Piotr Franciszek}{Gumulski}{1815}{Rudzienko} oraz
Marianny Chłopik\indeksosoby[Chłopik]{Marianna}{Gumulska}{1819}{Zamienie},
noszący nazwisko Gomulski, mieszkają do~dnia dzisiejszego w~Mińsku Mazowieckim
oraz jego okolicach, przede wszyskim:

\begin{itemize}
    \item w~Anielinie (obecnie części Mińska Mazowieckiego): potomkowie
    najmłodszego syna Piotra i~Marianny --
    Stanisława\indeksosoby{Stanisław}{Gomulski}{1863}{Anielina} (ur. 1863~r.)
    -- patrz: ryc. \ref{fig:wgomulski_1967}, ryc. \ref{fig:jgomulski_1996},
    \item w~Wólce Mińskiej: potomkowie najstarszego syna Piotra i~Marianny --
    Jana\indeksosoby{Jan}{Gomulski}{1836}{Zamienie} (ur. 1836~r.) -- patrz:
    ryc. \ref{fig:egomulski_1957}, ryc. \ref{fig:jgomulski_1932}, ryc.
    \ref{fig:mgomulski_1957},
    \item w~Karolinie: potomkowie najstarszego syna Jana Gomulskiego
    (ur. 1836~r.) -- Józefa\indeksosoby{Józef}{Gomulski}{1859}{Anielina}
    (ur. 1859~r.) -- patrz: ryc. \ref{fig:wgomulski_1957}, ryc.
    \ref{fig:agomulska_1935}, ryc. \ref{fig:fgomulski_1972},
  \end{itemize}

Świadczą o~tym liczne nagrobki na cmentarzu parafialnym parafii pod wezwaniem
Narodzenia Najświętszej Maryi Panny w~Mińsku Mazowieckim (ul. Kościelna~1,
05-300 Mińsk Mazowiecki), gdzie spoczywa wielu członków rodziny Gomulskich.

% Przednia okładka podrozdziału
\newpage
\includepdf{Karolina_mapa_fin.png}

\phantomsection
\addcontentsline{toc}{section}{Karolina: 1896~r. -- \biezacyrok}

\lettrine[lines=5, lhang=0.1]{K}{arolina} została założona około 1825~roku
w~odległości trzech kilometrów na~południe od~Mińska Mazowieckiego, na~terenie
ówczesnego Królestwa Polskiego. Wieś ta~należała na~początku XIX~wieku
do~parafii pod wezwaniem Narodzenia Najświętszej Maryi Panny w~Mińsku
Mazowieckim. Najstarsza znana wzmianka o~wsi Karolinia pochodzi z~1825~roku --
jest to~wpis w~aktach stanu cywilnego gminy Mińsk Mazowiecki, z~pierwszego
lipca ww. roku, dotyczący narodzin Jakuba Walesiaka syna Rocha i~Katarzyny
z~domu Wronikowskiej (patrz \ref{fig:karolina_1825}). Poza wspomnianą rodziną
Walesiaków, Karolinę na~początku XIX~wieku zamieszkiwały również między innymi
takie rodziny pochodzenia chłopskiego jak: Tarczyńscy, Kwiatkowscy, Rokiccy,
Nalazkowie czy Dobrzyńscy.

\begin{figure}[!ht]
    \vspace*{0.5cm}
    \centering \includegraphics[width=1.0\linewidth]{
        Karolina_1825.jpg}
    \captionsetup{format=hang}
    \caption{Pierwszy znany wpis dotyczący wsi Karolina pochodzący 
    z~1825~roku z~akt stanu cywilnego gminy Mińsk Mazowiecki 
    \cite{par_minsk1}.}
    \label{fig:karolina_1825}
\end{figure}

Według \emph{Tabelli miast, wsi, osad, Królestwa Polskiego: z~wyrażeniem ich
położenia i~ludności} przygotowanej przez Biuro Kommissyi Rządowey Spraw
Wewnętrznych i~Policyi Królestwa Polskiego \cite{biuro_1827} w~1827~roku
w~Karolinie było 10~domów, w~których mieszkało łącznie 49~osób -- patrz ryc.
\ref{fig:karolina_1827}.

\begin{figure}[!ht]
    \centering \includegraphics[width=0.83\linewidth]{
        Karolina_1827.jpg}
    \captionsetup{format=hang}
    \caption{Fragment \emph{Tabelli miast, wsi, osad, Królestwa Polskiego:
    z~wyrażeniem ich położenia i~ludności} zawierający wpis dotyczący wsi
    Karolina z 1827~roku \cite{biuro_1827}.}
    \label{fig:karolina_1827}
\end{figure}

\begin{figure}[!ht]
    \vspace*{0.5cm}
    \centering \includegraphics[width=0.83\linewidth]{
        1867_Anna_Dobrzyńska_akt_chrztu_parafia_Mińsk_Mazowiecki_wpis_249.jpg}
    \captionsetup{format=hang}
    \caption{Akt chrztu Anny Dobrzyńskiej -- par. Mińsk Mazowiecki
    1867~rok (249/1867) \cite{par_minsk2}.}
    \label{fig:agomulska_1867}
\end{figure}

Pierwszym znanym członkiem rodziny Gomulskich, który zamieszkał w~Karoline był
Józef Gomulski\indeksosoby{Józef}{Gomulski}{1859}{Anielina} (ur.~1859~r.~--
zm.~1929~r.), najstarszy syn
Jana\indeksosoby{Jan}{Gomulski}{1836}{Zamienie}
i~Agnieszki\indeksosoby[Piotrkowicz]{Agnieszka}{Gomulska}{1840}{Budy
Barcząckie} z~domu Piotrkowicz. Dnia
7~lutego 1887~roku Józef Gomulski ożenił się z~Anną
Dobrzyńską\indeksosoby[Dobrzyńska]{Anna}{Gomulska}{1867}{Karolina}
(ur.~1867~r.~-- zm.~1935~r.) panną ze~wsi Karolina, córką Adama i~Marianny
(z~domu Ryniewicz) Dobrzyńskich (patrz: ryc. \ref{fig:agomulska_1867} oraz
\ref{fig:jgomulski_1887}). Para po~ślubie zamieszkała w~rodzinnej wsi Józefa
-- Wólce Mińskiej, gdzie urodziła się im~piątka najstarszych dzieci:

\begin{itemize}
\item Julianna\indeksosoby{Julianna}{Gomulska}{1888}{Wólka Mińska}
(ur.~1888~r.~-- zm.~?) -- patrz: ryc. \ref{fig:jgomulska_1888},
\item Franciszek\indeksosoby{Franciszek}{Gomulski}{1889}{Wólka Mińska}
(ur.~1889~r.~-- zm.~1889~r.) -- patrz: ryc. \ref{fig:fgomulski_1889},
\item Józef\indeksosoby{Józef}{Gomulski}{1891}{Wólka Mińska}
(ur.~1891~r.~-- zm.~1891~r.) -- patrz: ryc. \ref{fig:jgomulski_1891},
\item Wincenty\indeksosoby{Wincenty \flagaPL}{Gomulski}{1892}{Wólka Mińska}
(ur.~1892~r.~-- zm.~1957~r.) -- patrz: ryc. \ref{fig:wingomulski_1892},
\item Marianna\indeksosoby{Marianna}{Gomulska}{1894}{Wólka Mińska}
(ur.~1894~r.~-- zm.~1958~r.) -- patrz: ryc. \ref{fig:mgomulska_1894}.
\end{itemize}

\begin{figure}[!ht]
    \vspace*{0.5cm}
    \centering \includegraphics[width=0.9\linewidth]{
        1887_Stanisław_Gomulski_Anna_Dobrzyńska_akt_ślubu_parafia_Mińsk_Mazowiecki_wpis_26.jpg}
    \captionsetup{format=hang}
    \caption{Akt ślubu Stanisława Gomulskiego i~Anny Dobrzyńskiej -- par. 
    Mińsk Mazowiecki 1887~rok (26/1887) \cite{par_minsk2}.}
    \label{fig:jgomulski_1887}
\end{figure}

\begin{figure}[!ht]
    \vspace*{0.5cm}
    \centering \includegraphics[width=0.87\linewidth]{
        1888_Julianna_Gomulska_akt_chrztu_parafia_Mińsk_Mazowiecki_wpis_76.jpg}
    \captionsetup{format=hang}
    \caption{Akt chrztu Julianny Gomulskiej -- par. Mińsk Mazowiecki
    1888~rok (76/1888) \cite{par_minsk2}.}
    \label{fig:jgomulska_1888}
\end{figure}

\begin{figure}[!ht]
    \vspace*{0.5cm}
    \centering \includegraphics[width=0.87\linewidth]{
        1889_Franciszek_Gomulski_akt_chrztu_parafia_Mińsk_Mazowiecki_wpis_360.jpg}
    \captionsetup{format=hang}
    \caption{Akt chrztu Franciszka Gomulskiego -- par. Mińsk Mazowiecki
    1889~rok (360/1889) \cite{par_minsk2}.}
    \label{fig:fgomulski_1889}
\end{figure}

\begin{figure}[!ht]
    \vspace*{0.5cm}
    \centering \includegraphics[width=0.77\linewidth]{
        1891_Józef_Gomulski_akt_chrztu_parafia_Mińsk_Mazowiecki_wpis_16.jpg}
    \captionsetup{format=hang}
    \caption{Akt chrztu Józefa Gomulskiego -- par. Mińsk Mazowiecki
    1891~rok (16/1891) \cite{par_minsk2}.}
    \label{fig:jgomulski_1891}
\end{figure}

\begin{figure}[!ht]
    \vspace*{0.5cm}
    \centering \includegraphics[width=0.77\linewidth]{
        1892_Wincenty_Gomulski_akt_chrztu_parafia_Mińsk_Mazowiecki_wpis_118.jpg}
    \captionsetup{format=hang}
    \caption{Akt chrztu Wincentego Gomulskiego -- par. Mińsk Mazowiecki
    1892~rok (118/1892) \cite{par_minsk2}.}
    \label{fig:wingomulski_1892}
\end{figure}

\begin{figure}[!ht]
    \vspace*{0.5cm}
    \centering \includegraphics[width=1.0\linewidth]{
        1894_Marianna_Gomulska_akt_chrztu_parafia_Mińsk_Mazowiecki_wpis_358.jpg}
    \captionsetup{format=hang}
    \caption{Akt chrztu Marianny Gomulskiej -- par. Mińsk Mazowiecki
    1894~rok (358/1894) \cite{par_minsk2}.}
    \label{fig:mgomulska_1894}
\end{figure}

Około 1896~roku Józef\indeksosoby{Józef}{Gomulski}{1859}{Anielina}
oraz Anna\indeksosoby[Dobrzyńska]{Anna}{Gomulska}{1867}{Karolina} Gomulscy,
wraz z~trójką\footnote{Dwóm najstarszym synom Józefa i~Anny: 
Franciszkowi\indeksosoby{Franciszek}{Gomulski}{1889}{Wólka Mińska}
oraz Józefowi\indeksosoby{Józef}{Gomulski}{1891}{Wólka Mińska} nie udało
przeżyć wieku niemowlęcego.} swoich dzieci:
Julianną\indeksosoby{Julianna}{Gomulska}{1888}{Wólka Mińska},
Wincentym\indeksosoby{Wincenty \flagaPL}{Gomulski}{1892}{Wólka Mińska} oraz
Marianną\indeksosoby{Marianna}{Gomulska}{1894}{Wólka Mińska}, przeprowadzili
się do~sąsiedniej wsi -- Karoliny, gdzie przyszła na~świat siódemka ich
najmłodszych dzieci:

\begin{itemize}
\item Feliks\indeksosoby{Feliks}{Gomulski}{1896}{Karolina} (ur.~1896~r.~--
zm.~1972~r.) -- patrz: ryc. \ref{fig:fgomulski_1896},
\item Helena\indeksosoby[Gomulska]{Helena}{Łapińska}{1899}{Karolina}
(ur.~1899~r.~-- zm.~?) -- patrz: ryc. \ref{fig:hgomulska_1898},
\item Zofia\indeksosoby[Gomulska]{Zofia}{Bartnicka}{1900}{Karolina}
(ur.~1900~r.~-- zm.~1970~r.) -- patrz: ryc. \ref{fig:zgomulska_1900},
\item Wacław\indeksosoby{Wacław}{Gomulski}{1904}{Karolina}
(ur.~1904~r.~-- zm.~?) -- patrz: ryc. \ref{fig:wgomulski_1904},
\item Stanisław\indeksosoby{Stanisław}{Gomulski}{1906}{Karolina}
(ur.~1906~r.~-- zm.~1957~r.) -- patrz: ryc. \ref{fig:sgomulski_1906},
\item Henryk\indeksosoby{Henryk}{Gomulski}{1908}{Karolina}
(ur.~1908~r.~-- zm.~1908~r.) -- patrz: ryc. \ref{fig:hgomulski_1908},
\item Zuzanna \indeksosoby{Zuzanna}{Gomulska}{1910}{Karolina}
(ur.~1910~r.~-- zm.~1911~r.) -- patrz: ryc. \ref{fig:zgomulska_1910}.
\end{itemize}

\begin{figure}[!ht]
    \vspace*{0.5cm}
    \centering \includegraphics[width=0.75\linewidth]{
        1896_Feliks_Gomulski_akt_chrztu_parafia_Mińsk_Mazowiecki_wpis_490.jpg}
    \captionsetup{format=hang}
    \caption{Akt chrztu Feliksa Gomulskiego -- par. Mińsk Mazowiecki
    1896~rok (490/1896) \cite{par_minsk2}.}
    \label{fig:fgomulski_1896}
\end{figure}

\begin{figure}[!ht]
    \vspace*{0.5cm}
    \centering \includegraphics[width=0.75\linewidth]{
        1898_Helena_Gomulska_akt_chrztu_parafia_Mińsk_Mazowiecki_wpis_549.jpg}
    \captionsetup{format=hang}
    \caption{Akt chrztu Heleny Gomulskiej -- par. Mińsk Mazowiecki
    1898~rok (549/1898) \cite{par_minsk2}.}
    \label{fig:hgomulska_1898}
\end{figure}

\begin{figure}[!ht]
    \vspace*{0.5cm}
    \centering \includegraphics[width=0.73\linewidth]{
        1900_Zofia_Gomulska_akt_chrztu_parafia_Mińsk_Mazowiecki_wpis_527.jpg}
    \captionsetup{format=hang}
    \caption{Akt chrztu Zofii Gomulskiej -- par. Mińsk Mazowiecki
    1900~rok (527/1900) \cite{par_minsk2}.}
    \label{fig:zgomulska_1900}
\end{figure}

\begin{figure}[!ht]
    \vspace*{0.5cm}
    \centering \includegraphics[width=0.73\linewidth]{
        1904_Wacław_Gomulski_akt_chrztu_parafia_Mińsk_Mazowiecki_wpis_58.jpg}
    \captionsetup{format=hang}
    \caption{Akt chrztu Wacława Gomulskiego -- par. Mińsk Mazowiecki
    1904~rok (58/1904) \cite{par_minsk2}.}
    \label{fig:wgomulski_1904}
\end{figure}

\begin{figure}[!ht]
    \vspace*{0.5cm}
    \centering \includegraphics[width=0.77\linewidth]{
        1906_Stanisław_Gomulski_akt_chrztu_parafia_Mińsk_Mazowiecki_wpis_658.jpg}
    \captionsetup{format=hang}
    \caption{Akt chrztu Stanisława Gomulskiego -- par. Mińsk Mazowiecki
    1906~rok (658/1906) \cite{par_minsk2}.}
    \label{fig:sgomulski_1906}
\end{figure}

\begin{figure}[!ht]
    \vspace*{0.5cm}
    \centering \includegraphics[width=0.77\linewidth]{
        1908_Henryk_Gomulski_akt_chrztu_parafia_Mińsk_Mazowiecki_wpis_484.jpg}
    \captionsetup{format=hang}
    \caption{Akt chrztu Henryka Gomulskiego -- par. Mińsk Mazowiecki
    1908~rok (484/1908) \cite{par_minsk2}.}
    \label{fig:hgomulski_1908}
\end{figure}

\begin{figure}[!ht]
    \vspace*{0.5cm}
    \centering \includegraphics[width=1.0\linewidth]{
        1910_Zuzanna_Gomulska_akt_chrztu_parafia_Mińsk_Mazowiecki_wpis_463.jpg}
    \captionsetup{format=hang}
    \caption{Akt chrztu Zuzanny Gomulskiej -- par. Mińsk Mazowiecki
    1910~rok (463/1910) \cite{par_minsk2}.}
    \label{fig:zgomulska_1910}
\end{figure}

\begin{figure}[!ht]
    \vspace*{0.5cm}
    \centering \includegraphics[width=1.0\linewidth]{
        Wincenty_Gomulski_Medal_Niepodległości_1.jpg}
    \captionsetup{format=hang}
    \caption{Fragment wniosku o~przyznanie Medalu Niepodległości Wincentowi
    Gomulskiemu znajdujący się w~archiwum Wojskowego Biura Historycznego
    w~Warszawie pod sygnaturą \enquote{AP 147 podofic.} \cite{wgomulski}.}
    \label{fig:wgomulski_1938_1}
\end{figure}

\begin{figure}[!ht]
    \vspace*{0.5cm}
    \centering \includegraphics[width=0.84\linewidth]{
        Wincenty_Gomulski_Medal_Niepodległości_2.jpg}
    \captionsetup{format=hang}
    \caption{Fragment wniosku o~przyznanie Medalu Niepodległości Wincentowi
    Gomulskiemu znajdujący się w~archiwum Wojskowego Biura Historycznego
    w~Warszawie pod sygnaturą \enquote{AP 147 podofic.} \cite{wgomulski}.}
    \label{fig:wgomulski_1938_2}
\end{figure}

\begin{figure}[!ht]
    \vspace*{0.5cm}
    \centering \includegraphics[width=0.84\linewidth]{
        Wincenty_Gomulski_Medal_Niepodległości_3.jpg}
    \captionsetup{format=hang}
    \caption{Fragment wniosku o~przyznanie Medalu Niepodległości Wincentowi
    Gomulskiemu znajdujący się w~archiwum Wojskowego Biura Historycznego
    w~Warszawie pod sygnaturą \enquote{AP 147 podofic.} \cite{wgomulski}.}
    \label{fig:wgomulski_1938_3}
\end{figure}

Spośród dwanaściorga dzieci Józefa\indeksosoby{Józef}{Gomulski}{1859}{Anielina}
i~Anny\indeksosoby[Dobrzyńska]{Anna}{Gomulska}{1867}{Karolina} Gomulskich
prawdopodobnie tylko szóstka dożyła wieku dorosłego:
Wincenty\indeksosoby{Wincenty \flagaPL}{Gomulski}{1892}{Wólka Mińska},
Marianna\indeksosoby{Marianna}{Gomulska}{1894}{Wólka Mińska},
Feliks\indeksosoby{Feliks}{Gomulski}{1896}{Karolina},
Helena\indeksosoby{Helena}{Łapińska}{1899}{Karolina},
Zofia\indeksosoby{Zofia}{Bartnicka}{1900}{Karolina}
i~Stanisław\indeksosoby{Stanisław}{Gomulski}{1906}{Karolina}.
Poza Wincentym, wszyscy osiedlili się wraz z~rodzinami
na~stałe w~Karolinie. Wincenty
Gomulski\indeksosoby{Wincenty \flagaPL}{Gomulski}{1892}{Wólka Mińska}
natomiast przeprowadził się do~miejscowości Postawy znajdującej się na~terenie
obecnej Białorusi. Świadczy o~tym wniosek z~1938~roku o~przyznanie mu~Medalu
Niepodległości za~udział w~walkach o~niepodległość Polski w~latach 1918--1921,
gdzie wskazano, iż~w chwili składania wniosku mieszkał on~w~miejscowości
Postawy przy placu imienia marszałka Józefa Piłsudskiego (patrz: ryc.
\ref{fig:wgomulski_1938_1}, \ref{fig:wgomulski_1938_2},
\ref{fig:wgomulski_1938_3}). Wincenty
Gomulski\indeksosoby{Wincenty \flagaPL}{Gomulski}{1892}{Wólka Mińska} w~Wojsku
Polskim dosłużył się stopnia wachmistrza (starszego podoficera). Informacja
o~odznaczeniu Wincentego Medalem Niepodległości została opublikowana
w~1938~roku w~\emph{Monitorze Polskim} \cite{mp_1938}. Można ją również
odnaleźć na~stronie internetowej Fundacji Obchodów 100.~Rocznicy Odzysknia
Przez Polskę
Niepodległości\footnote{\url{https://fundacja100.pl/krzyz-i-medal-niepodleglosci/lista-odznaczonych/wincenty-gomulski}} --
patrz: ryc. \ref{fig:wgomulski_fundacja}.

\begin{figure}[!ht]
    \vspace*{0.5cm}
    \centering \includegraphics[width=0.95\linewidth]{
        Wincenty_Gomulski_bojownik_niepodległości.png}
    \captionsetup{format=hang}
    \caption{Wpis dotyczący Wincentego Gomulskiego pochodzący ze~strony
    internetowej Fundacji Obchodów 100.~Rocznicy Odzysknia Przez Polskę
    Niepodległości.}
    \label{fig:wgomulski_fundacja}
\end{figure}

Po drugiej wojnie światowej Wincenty
Gomulski\indeksosoby{Wincenty \flagaPL}{Gomulski}{1892}{Wólka Mińska}
prawdopodobnie powrócił do~rodzinnej wsi Karolina. Wskazuje na~to~notatka
na~marginesie jego aktu chrztu (patrz ryc. \ref{fig:wingomulski_1892}),
która mówi, iż~dnia 26~czerwca 1951~roku, w~wieku 59~lat, wziął on~ślub
w~urzędzie stanu cywilnego w~Rembertowie z~Antoniną Marsz\footnote{Był to~jego
drugi ślub, bo~w~1919~roku w~Warszawie ożenił się z~Janiną Panek.}. Dodatkowo
jego powrót w~rodzinne strony uwiarygodnia fakt, iż~po~śmierci 7~lutego
1957~roku, został on~pochowany na~cmentarzu parafialnym parafii pod wezwaniem
Narodzenia Najświętszej Maryi Panny w~Mińsku Mazowieckim -- patrz ryc.
\ref{fig:wgomulski_1957}.

\begin{figure}[!ht]
    \vspace*{0.5cm}
    \centering \includegraphics[width=0.75\linewidth]{
        1957_Wincenty_Gomulski_zdjęcie_nagrobka1_cmentarz_Mińsk_Mazowiecki.jpg}
    \captionsetup{format=hang}
    \caption{Zdjęcie grobu Wincentego Gomulskiego, znajdującego się na
    cmentarzu parafialnym parafii pod wezwaniem Narodzenia Najświętszej Maryi
    Panny w~Mińsku Mazowieckim (sektor: IX, rząd: 12, numer: 10)}
    \label{fig:wgomulski_1957}
\end{figure}

\begin{figure}[!ht]
    \vspace*{0.5cm}
    \centering \includegraphics[width=0.75\linewidth]{
        1935_Anna_Gomulska_Dobrzyńska_zdjęcie_nagrobka1_cmentarz_Mińsk_Mazowiecki.jpg}
    \captionsetup{format=hang}
    \caption{Zdjęcie grobu Anny Gomulskiej z~d.~Dobrzyńskiej oraz
    Stanisława Gomulskiego, znajdującego się na cmentarzu parafialnym parafii
    pod wezwaniem Narodzenia Najświętszej Maryi Panny w~Mińsku Mazowieckim
    (sektor: XVIII, rząd: 1, numer: 6)}
    \label{fig:agomulska_1935}
\end{figure}

\begin{figure}[!ht]
    \vspace*{0.4cm}
    \centering \includegraphics[width=0.77\linewidth]{
        1936_Anna_Gomulska_Bartnicka_zdjęcie_nagrobka1_cmentarz_Mińsk_Mazowiecki.jpg}
    \captionsetup{format=hang}
    \caption{Zdjęcie grobu Anny Gomulskiej, znajdującego się na~cmentarzu
    parafialnym parafii pod wezwaniem Narodzenia Najświętszej Maryi Panny
    w~Mińsku Mazowieckim (sektor: XVIII, rząd: 1, numer: 45)}
    \label{fig:agomulska_1936}
\end{figure}

\begin{figure}[!ht]
    \vspace*{0.4cm}
    \centering \includegraphics[width=0.77\linewidth]{
        1970_Zofia_Gomulska_Bartnicka_zdjęcie_nagrobka1_cmentarz_Mińsk_Mazowiecki.jpg}
    \captionsetup{format=hang}
    \caption{Zdjęcie grobu małżonków Zofii oraz Władysława Bartnickich,
    znajdującego się na~cmentarzu parafialnym parafii pod wezwaniem Narodzenia
    Najświętszej Maryi Panny w~Mińsku Mazowieckim (sektor: IX, rząd: 10,
    numer: 6)}
    \label{fig:zgomulska_1970}
\end{figure}

\begin{figure}[!ht]
    \vspace*{0.4cm}
    \centering \includegraphics[width=0.77\linewidth]{
        1972_Feliks_Gomulski_zdjęcie_nagrobka1_cmentarz_Mińsk_Mazowiecki.jpg}
    \captionsetup{format=hang}
    \caption{Zdjęcie grobu Marianny i~Feliksa Gomulskich, znajdującego się na
    cmentarzu parafialnym parafii pod wezwaniem Narodzenia Najświętszej Maryi
    Panny w~Mińsku Mazowieckim (sektor: VIII, rząd: 4, numer: 3)}
    \label{fig:fgomulski_1972}
\end{figure}

Bezpośredni potomkowie
Józefa\indeksosoby{Józef}{Gomulski}{1859}{Anielina}
i~Anny\indeksosoby[Dobrzyńska]{Anna}{Gomulska}{1867}{Karolina} Gomulskich,
mieszkają na~terenie obecnej Karoliny do~dnia dzisiejszego -- ich groby można
odnaleźć na~cmentarzu parafialnym parafii pod wezwaniem Narodzenia
Najświętszej Maryi Panny w~Mińsku Mazowieckim. Ryciny: \ref{fig:wgomulski_1957},
\ref{fig:agomulska_1935}, \ref{fig:agomulska_1936}, \ref{fig:zgomulska_1970} 
oraz \ref{fig:fgomulski_1972} przedstawiają niektóre z~tych nagrobków.